% !TeX root = ../../Book4.tex
% 纲要标题：## 第8章　Moser 方法、Darboux 与邻域定理
% 纲要位置：几何学计划纲要.md，第 3775 行。
% 正文状态：初稿；尚未正式审核。

% 纲要内容开始
% 
% > 本章目标与证明范围：同一变形方法依次处理整体同位、局部标准形与沿子流形的邻域。每个版本均说明原始一形式、非退化区间及向量场完备性。
% 
% 纲要内容结束

\chapter{Moser 方法、Darboux 与邻域定理}
\label{g4:ch:08}
\BookChapterNumber{g4:ch:08}

\begin{chapterabstract}
本章用时间依赖向量场比较一族辛形式。拉回求导把几何问题化为一个逐点线性方程：选取变化率的原始一形式，再用辛形式的非退化性求出向量场。在闭流形上，这给出上同调类不变的辛形式族的同位；在一点附近，径向同伦与流的时间估计给出 Darboux 坐标；沿子流形，匹配环境切空间上的线性数据后，得到辛子流形与 Lagrange 子流形的邻域定理。最后把辛同胚的图像写成 Lagrange 子流形，引入局部生成函数，并以非挤压现象说明局部标准形之外仍有整体限制。
\end{chapterabstract}

\begin{learninggoals}
\begin{itemize}
	\item 推导随时间变化的拉回公式，构造光滑原始一形式族，并证明闭流形上的 Moser 稳定性定理。
	\item 分别核查原始形式、整条路径的非退化性和流的存在时间，识别非紧情形中缺少的条件。
	\item 用径向同伦和局部流证明带指定辛标架的 Darboux 定理，理解辛形式没有局部曲率不变量的含义。
	\item 区分向子流形拉回为零与整个环境张量为零，证明保留一阶数据的相对 Moser 定理。
	\item 由辛法丛及余切丛构造邻域模型，说明辛子流形和 Lagrange 子流形各需要哪些数据。
	\item 用图像判别辛同胚，写出两类生成函数及其投影条件，并区分面积、体积与辛嵌入的限制。
\end{itemize}
\end{learninggoals}

上一章已经把每根切空间上的辛形式化为标准配对。然而，逐点选出的辛基未必来自一组坐标的坐标向量场。要得到实际坐标，还要利用二形式的闭性。本章让坐标随时间运动：如果一族形式的变化恰好被坐标运动抵消，那么移动后的拉回形式就保持不变。

\ProseParagraphBreak

这个思路有三个需要分别解决的环节。变化率必须有光滑选取的原始一形式；中间各时刻的二形式必须非退化；得到的向量场必须能够运行到指定终点。Moser 方法把这三件事放在同一个证明里，但它们的依据不同：上同调或同伦公式提供原始，非退化性的开性控制小邻域，紧致性或轨道估计控制存在时间。把这些依据说清楚，同一方法便能从整体同位转到一点及子流形附近。

\ProseParagraphBreak

本章使用第一卷的逆函数定理、光滑流与微分形式，以及第二卷的 de Rham 上同调和管状邻域。光滑原始族的证明还用到第三卷定理2.2.6所给的小强凸邻域，以便选取有限好覆盖；这里只使用该结论构造固定的局部同伦。邻域定理中的线性配对、相容结构及余切丛符号都沿用上一章，特别是典范辛形式取为 \(-\mathrm{d}\theta\)。

\ProseParagraphBreak

全章的证明主线到 Lagrange 邻域定理为止是完整的。图像与生成函数随后给出具体应用；末节陈述的非挤压定理将在\cref{g4:ch:24}证明，本章先证明其线性版本，并与保体积变换作比较。有关变形方法、邻域模型与生成函数，可对照 Cannas da Silva 正式综述的第1.4—1.6及2.2—2.5节。\cite[第1.4--1.6、2.2--2.5节]{Cannas2006Symplectic}

% !TeX root = ../../Book4.tex
% 纲要标题：### 8.1 时间依赖流与 Moser 方程
% 纲要位置：几何学计划纲要.md，第 3779 行。
% 正文状态：初稿；尚未正式审核。

% 纲要内容开始
% 
% * Lie 导数和拉回求导
% * 同上同调类的形式族
% * 求解收缩方程
% * 紧性、紧支和相对保持
% 
% 纲要内容结束

\section{时间依赖流与 Moser 方程}
\label{g4:sec:0801}

设一族辛形式 \(\omega_t\) 随时间变化。我们希望同时移动流形上的点，使移动后的配对不再变化，即找到 \(\phi_t\) 满足 \(\phi_t^*\omega_t=\omega_0\)。直接求这样一族微分同胚并不容易，但对时间求导后，问题会变成逐点的线性方程。这个思路需要分别解决三个问题：形式自身的变化怎样与点的移动相加，恰当形式的原始形式怎样随参数光滑选择，以及所得向量场能否流动到指定时间。

\subsection{Lie 导数和拉回求导}
\label{g4:subsec:0801:01}

本节的流形均为有限维、Hausdorff、第二可数的无边界光滑流形。区间 \([0,1]\) 上的光滑族是指系数对时间和空间变量联合光滑，在时间端点按局部光滑延拓理解。记 \(\dot\eta_t=\partial_t\eta_t\)，而 \(\mathrm{d}\eta_t\) 始终表示对空间变量作外微分。

\begin{definition}\label{g4:def:0801:time-dependent-flow}
	光滑映射 \(X:[0,1]\times M\to\mathrm{T}M\)，满足 \(X_t(p)=X(t,p)\in\mathrm{T}_pM\)，称为\term[shijianyilaixiangliangchang]{时间依赖向量场}[time-dependent vector field]。其\term[yanhuayingshe]{演化映射}[evolution map] \(\Phi_{t,s}\) 由初值问题
	\[
		\partial_t\Phi_{t,s}(p)=X_t(\Phi_{t,s}(p)),
		\qquad \Phi_{s,s}(p)=p
	\]
	确定，并只在相应解存在时定义。若 \(\phi_t:M\to M\) 是光滑微分同胚族且 \(\phi_0=\operatorname{id}\)，则称它为从恒等映射出发的同位\termidx[tongwei]{同位}。这沿用第6章对微分同胚之间同位的定义。
\end{definition}

第一卷命题13.1.9证明了演化映射的存在、唯一、联合光滑性和定义域开放性。其证明将方程化为乘积流形上的自主向量场 \(\partial_t+X_t\)，因而时间端点的局部延拓也能用同一方法处理。演化映射满足
\begin{equation}\label{g4:eq:0801:evolution-law}
	\Phi_{t,r}\circ\Phi_{r,s}=\Phi_{t,s},\qquad
	\Phi_{t,s}^{-1}=\Phi_{s,t},
\end{equation}
其中等式在相应开定义域上成立。以下常写 \(\phi_t=\Phi_{t,0}\)。在非紧流形上，\(\phi_t\) 可能只在某个开集 \(U\subset M\) 上存在；此时它仍是 \(U\) 到其开像的微分同胚。

\ProseParagraphBreak

对一个不随时间变化的向量场 \(Y\)，沿其正向流拉回求导定义 Lie 导数 \(\mathcal L_Y\)。第一卷定理14.1.5给出的 Cartan 公式\footnote{Élie Joseph Cartan（1869--1951），法国数学家，研究 Lie 群、微分几何与外微分系统。}为
\begin{equation}\label{g4:eq:0801:cartan}
	\mathcal L_Y\eta=\mathrm{d}(\iota_Y\eta)+\iota_Y(\mathrm{d}\eta).
\end{equation}
这里插入放在第一个变量，即 \((\iota_Y\eta)(v_1,\ldots)=\eta(Y,v_1,\ldots)\)。对 \(X_t\) 使用 \(\mathcal L_{X_t}\) 时，先固定时刻 \(t\)，再把 \(X_t\) 作为一个普通向量场。

\begin{proposition}\label{g4:prop:0801:pullback-derivative}
	设 \(M\) 为无边界光滑流形，\(X_t\) 是光滑时间依赖向量场，\(\eta_t\) 是光滑 \(k\) 形式族。若其从时刻 \(0\) 出发的流 \(\phi_t\) 在 \([0,1]\times U\) 上存在，其中 \(U\subset M\) 开，则
	\begin{equation}\label{g4:eq:0801:pullback-derivative}
		\frac{\mathrm{d}}{\mathrm{d}t}(\phi_t^*\eta_t)
		=\phi_t^*\bigl(\dot\eta_t+\mathcal L_{X_t}\eta_t\bigr).
	\end{equation}
	同一公式也适用于任意光滑微分同胚族，只需令 \(X_t=\dot\phi_t\circ\phi_t^{-1}\)，并在其像上解释右端。
\end{proposition}
\begin{proof}
	固定时刻 \(t\)，令 \(\psi_h=\Phi_{t+h,t}\)。由演化复合律，局部有 \(\phi_{t+h}=\psi_h\circ\phi_t\)，所以
	\[
		\phi_{t+h}^*\eta_{t+h}=\phi_t^*(\psi_h^*\eta_{t+h}).
	\]
	\(\psi_0\) 为恒等映射，且 \(\partial_h\psi_h|_{h=0}=X_t\)。因此对固定函数 \(f\)，有 \(\partial_h(\psi_h^*f)|_0=X_tf\)；对函数微分有
	\[
		\left.\partial_h\right|_0\psi_h^*(\mathrm{d}f)
		=\mathrm{d}(X_tf).
	\]
	后一式使用第一卷命题14.1.8的拉回与外微分相容性，以及光滑系数的时间微分与空间微分可交换。任意形式局部是函数系数与坐标微分的楔积之和，由乘积求导法则，\(\partial_h(\psi_h^*\eta)|_0=\mathcal L_{X_t}\eta\)。形式本身的时间变化再贡献 \(\dot\eta_t\)，故
	\[
		\left.\partial_h\right|_0(\psi_h^*\eta_{t+h})
		=\dot\eta_t+\mathcal L_{X_t}\eta_t.
	\]
	拉回到 \(U\) 即得结论。对一般微分同胚族，改取 \(\psi_h=\phi_{t+h}\circ\phi_t^{-1}\)，同一论证成立。上述局部等式在各坐标域相容，并由光滑性延至时间端点。
\end{proof}

公式中的两项分别来自形式的变化和输入点的移动，二者都不可省略。特别地，若 \(\mathrm{d}\omega_t=0\)，则
\begin{equation}\label{g4:eq:0801:closed-pullback-derivative}
	\frac{\mathrm{d}}{\mathrm{d}t}(\phi_t^*\omega_t)
	=\phi_t^*\bigl(\dot\omega_t+\mathrm{d}(\iota_{X_t}\omega_t)\bigr).
\end{equation}
我们将让括号内的两个二形式彼此抵消。

\subsection{同上同调类的形式族}
\label{g4:subsec:0801:02}

先明确上同调条件的含义。对闭二形式，\([\omega_t]=[\omega_0]\) 在 \(H^2_{\mathrm{dR}}(M)\) 中成立，按定义等价于 \(\omega_t-\omega_0\) 恰当。第二卷定理7.3.11把这一群与实奇异上同调自然识别，所以也可用拓扑上同调表达同一条件。逐个时刻选择一个原始一形式还不能直接用于微分方程，因为这些选择可能彼此无关。下面将局部到整体的构造固定下来，使原始形式随数据光滑变化。

\begin{lemma}\label{g4:lem:0801:smooth-primitives}
	设 \(M\) 是紧无边界光滑流形。可以为其所有恰当二形式选择一个实线性算子
	\[
		T:\mathrm{d}\Omega^1(M)\longrightarrow\Omega^1(M),
		\qquad \mathrm{d}(T\eta)=\eta,
	\]
	使得任意光滑参数族 \(\eta_t\in\mathrm{d}\Omega^1(M)\) 都给出光滑族 \(T\eta_t\)。构造与参数的任意阶微分交换，也适用于有限个光滑参数。
\end{lemma}
\begin{proof}
	我们先在有限覆盖的各片取原始形式，再修正交叠上的差。唯一需要整体求解的部分将落在一个有限维线性方程中；固定它的一个线性右逆，就能避免逐个时刻作不相容的选择。

	先固定一个有限光滑好覆盖 \((U_i)_{i=1}^N\)，即每个非空有限交都能在自身内光滑收缩到一点。这里说明其存在，以便后面的收缩算子确有光滑性。第一卷定理12.3.11给出一个辅助黎曼度量；第三卷定理2.2.6给出每点附近的强凸开邻域，其中任意两点由唯一最短测地线连接，且该线光滑依赖端点并留在邻域中。由紧性选有限个这样的邻域覆盖 \(M\)。若有限交非空，其中两点在各邻域里的最短线由整个流形中的唯一性而相同，故整条线留在交集中。固定交集中的一点，沿这些最短线就得到光滑收缩。这正是第二卷定义7.3.9所用的光滑好覆盖；不连通的 \(M\) 可逐个分支作上述构造。

	对每个需要使用的非空交集 \(U_I\)，固定基点 \(a_I\) 及光滑同伦 \(H_I:[0,1]\times U_I\to U_I\)，使 \(H_I(0,x)=a_I\)、\(H_I(1,x)=x\)。第一卷定理14.2.1的同伦算子给出固定线性算子 \(K_I\)，满足
	\[
		\mathrm{d}K_I\xi+K_I\mathrm{d}\xi=\xi
		\quad\text{对正次数形式 }\xi.
	\]
	其具体定义为沿 \([0,1]\) 积分 \(H_I^*\xi\) 的时间插入分量，因而将光滑参数族送到光滑参数族。任意阶参数微分都可移入该紧区间积分。

	现在取一个恰当二形式 \(\eta\)。令
	\[
		\beta_i=K_i(\eta|_{U_i}),\qquad \mathrm{d}\beta_i=\eta|_{U_i}.
	\]
	在非空 \(U_{ij}\) 上，\(\beta_j-\beta_i\) 闭。对 \(i<j\) 令
	\[
		f_{ij}=K_{ij}(\beta_j-\beta_i),\qquad
		\mathrm{d}f_{ij}=\beta_j-\beta_i,
	\]
	再规定 \(f_{ji}=-f_{ij}\)、\(f_{ii}=0\)。在三重交上，函数
	\[
		c_{ijk}=f_{jk}-f_{ik}+f_{ij}
	\]
	的外微分为零。交集连通，所以它是一个实常数；可以通过在已固定的基点取值来确定这个常数。

	把每个非空 \((r+1)\) 重交对应的一份 \(\mathbb R\) 相乘，得到有限维常系数 \term[cechshanglian]{Čech 上链}[Čech cochain]空间 \(C^r\)。指标按交替规则排列，其微分为
	\[
		(\delta a)_{i_0\cdots i_{r+1}}
		=\sum_{\nu=0}^{r+1}(-1)^\nu
		a_{i_0\cdots\widehat{i_\nu}\cdots i_{r+1}}.
	\]
	刚才的常数构成 \(c\in C^2\)，且 \(\delta c=0\)。我们还需要它为余边界，这里直接使用 \(\eta\) 的恰当性验证。暂取任意一个整体原始形式 \(\zeta\)，满足 \(\mathrm{d}\zeta=\eta\)。在 \(U_i\) 上，\(\beta_i-\zeta\) 闭，故可写成 \(\mathrm{d}h_i\)。于是
	\[
		b_{ij}=f_{ij}-h_j+h_i
	\]
	在 \(U_{ij}\) 上为常数，且 \(\delta b=c\)。因此 \(c\in\operatorname{im}(\delta:C^1\to C^2)\)。这个临时的 \(\zeta\) 只用于证明可解性，不参与最终选择。

	有限维线性代数允许预先固定一个右逆
	\[
		R:\operatorname{im}(\delta:C^1\to C^2)\longrightarrow C^1,
		\qquad \delta R=\operatorname{id}.
	\]
	例如选 \(\ker\delta\) 的线性补空间，限制 \(\delta\) 后取逆即可。令 \(a=Rc\)、\(\widetilde f_{ij}=f_{ij}-a_{ij}\)，则
	\[
		\widetilde f_{jk}-\widetilde f_{ik}+\widetilde f_{ij}=0,
		\qquad \mathrm{d}\widetilde f_{ij}=\beta_j-\beta_i.
	\]
	取从属于 \((U_i)\) 的固定光滑单位分解 \((\rho_i)\)，其存在由第一卷定理12.3.5保证。定义 \(U_i\) 上的函数
	\[
		g_i=\sum_k\rho_k\widetilde f_{ki}.
	\]
	每项先在 \(U_k\cap U_i\) 定义，再向 \(U_i\) 的其余部分补零。因为 \(\operatorname{supp}\rho_k\subset U_k\)，补零在每点附近光滑；这里总和有限。三重交上的等式给出
	\[
		g_j-g_i=\sum_k\rho_k(\widetilde f_{kj}-\widetilde f_{ki})
		=\sum_k\rho_k\widetilde f_{ij}=\widetilde f_{ij}.
	\]
	因此 \(\beta_i-\mathrm{d}g_i\) 在交集上相等，粘成一个整体一形式，定义为 \(T\eta\)。它满足 \(\mathrm{d}(T\eta)=\eta\)。这也是第二卷引理7.3.10中单位分解收缩在当前次数的具体形式。

	所有实际操作都是固定收缩下的参数积分、限制、在固定点取值、固定有限维线性映射、乘固定光滑函数、有限求和以及外微分。这些操作都线性，并将光滑族送到光滑族，与任意阶参数微分交换。求 \(R\) 时使用的有限维子空间是闭子空间，故 \(c_t\) 的各阶参数导数仍在其定义域中。对 \(\mathrm{d}(T\eta_t)=\eta_t\) 求导还表明这些参数导数仍为恰当形式，所以与 \(T\) 交换的表述也有定义。于是 \(T\) 具有所述全部性质。
\end{proof}

\begin{corollary}\label{g4:cor:0801:constant-class-derivative}
	设 \(M\) 紧且无边界，\(\omega_t\) 为光滑闭二形式族，且 \([\omega_t]=[\omega_0]\) 对 \(0\le t\le1\) 成立。则存在光滑一形式族 \(\alpha_t\)，满足
	\begin{equation}\label{g4:eq:0801:primitive-derivative}
		\dot\omega_t=\mathrm{d}\alpha_t.
	\end{equation}
\end{corollary}
\begin{proof}
	差值 \(\eta_t=\omega_t-\omega_0\) 是光滑恰当族。用上一引理令 \(\kappa_t=T\eta_t\)，则 \(\mathrm{d}\kappa_t=\omega_t-\omega_0\)，且 \(\kappa_t\) 联合光滑。现在取 \(\alpha_t=\partial_t\kappa_t\)，交换系数的时间与空间微分，就得到 \(\mathrm{d}\alpha_t=\dot\omega_t\)。
\end{proof}

这说明上同调假设如何转化为一个可用于常微分方程的光滑右端。若一开始已有 \(\omega_t=\omega_0+\mathrm{d}\sigma_t\)，其中 \(\sigma_t\) 光滑，就直接取 \(\alpha_t=\dot\sigma_t\)。尤其对直线路径 \(\omega_t=\omega_0+t\,\mathrm{d}\sigma\)，可取不随时间变化的 \(\alpha_t=\sigma\)。

\subsection{求解收缩方程}
\label{g4:subsec:0801:03}

设 \(\omega_t\) 是辛形式族，并已找到满足式\eqref{g4:eq:0801:primitive-derivative}的 \(\alpha_t\)。式\eqref{g4:eq:0801:closed-pullback-derivative}提示我们逐点求解
\begin{equation}\label{g4:eq:0801:moser-equation}
	\iota_{X_t}\omega_t=-\alpha_t.
\end{equation}
这个等式称为\term[moserfangcheng]{Moser 方程}[Moser equation]。\footnote{Jürgen Kurt Moser（1928--1999），德国出生的美籍数学家，研究 Hamilton 动力系统与偏微分方程。}负号使 \(\mathrm{d}(\iota_{X_t}\omega_t)\) 与 \(\dot\omega_t\) 抵消。

\begin{lemma}\label{g4:lem:0801:contraction-solution}
	设 \(M\) 为光滑流形，\(\omega_t\) 为光滑非退化二形式族，\(\alpha_t\) 为光滑一形式族。方程\eqref{g4:eq:0801:moser-equation}有唯一光滑时间依赖向量场解。对任意 \((t,p)\)，有 \(X_t(p)=0\) 当且仅当 \(\alpha_t(p)=0\) 作为 \(\mathrm{T}_p^*M\) 中的整个协向量成立。
\end{lemma}
\begin{proof}
	非退化性说明 \(\omega_t^\flat:\mathrm{T}_pM\to\mathrm{T}_p^*M\)，\(v\mapsto\omega_t(v,\cdot)\)，是线性同构，所以唯一解为 \(X_t=-(\omega_t^\flat)^{-1}\alpha_t\)。在任一局部标架中，这个逆的矩阵由伴随矩阵除以非零行列式给出，因此随 \((t,p)\) 光滑。各坐标解由唯一性一致，构成整体光滑场。线性同构把零与零对应，给出最后的断言。
\end{proof}

例如，在 \(\omega=\sum_i\mathrm{d}q^i\wedge\mathrm{d}p_i\) 的坐标中，若 \(\alpha=\sum_i(a_i\,\mathrm{d}q^i+b_i\,\mathrm{d}p_i)\)，则方程的解为
\[
	X=\sum_i\left(-b_i\partial_{q^i}+a_i\partial_{p_i}\right).
\]
这直接核对了收缩所用的符号。原始形式的选择一般不唯一；若把 \(\alpha_t\) 改为 \(\alpha_t+\lambda_t\)，其中 \(\lambda_t\) 闭，所得速度场也会改变，但仍使式\eqref{g4:eq:0801:closed-pullback-derivative}的右端为零。

\begin{theorem}[Moser 稳定性定理]\label{g4:thm:0801:moser}
	设 \(M\) 为紧无边界光滑流形，\(\omega_t\)、\(0\le t\le1\)，为光滑辛形式族。若其 de Rham 上同调类 \([\omega_t]\) 与 \(t\) 无关，则存在从恒等映射出发的光滑同位 \(\phi_t:M\to M\)，使
	\begin{equation}\label{g4:eq:0801:moser-pullback}
		\phi_t^*\omega_t=\omega_0\qquad(0\le t\le1).
	\end{equation}
	反过来，满足该等式的同位必使 \([\omega_t]\) 与 \(t\) 无关。
\end{theorem}
\begin{proof}
	由\cref{g4:cor:0801:constant-class-derivative}选择光滑 \(\alpha_t\) 使 \(\dot\omega_t=\mathrm{d}\alpha_t\)，再由\cref{g4:lem:0801:contraction-solution}解出 \(X_t\)。紧性在这里保证流能够走完整个时间区间。具体而言，第一卷命题13.2.10证明共同紧支的时间依赖场有整体演化映射；现在可取共同紧集为整个 \(M\)。该命题的延拓论证应用于增广场 \(\partial_t+X_t\)：若解在有限时刻终止，轨道却一直位于紧的时间段与 \(M\) 的乘积中，就与有限端点逃逸判据矛盾。端点处用局部光滑延拓即可。

	于是 \(\phi_t=\Phi_{t,0}\) 对所有初值及 \(0\le t\le1\) 存在。反向演化给出光滑逆，故它是从恒等映射出发的同位。由拉回求导公式及 \(\mathrm{d}\omega_t=0\)，
	\[
		\frac{\mathrm{d}}{\mathrm{d}t}(\phi_t^*\omega_t)
		=\phi_t^*\mathrm{d}(\alpha_t+\iota_{X_t}\omega_t)=0.
	\]
	初值为 \(\omega_0\)，得到式\eqref{g4:eq:0801:moser-pullback}。

	反过来，设这样的同位已经存在，令 \(X_t=\dot\phi_t\circ\phi_t^{-1}\)。对等式求导并利用拉回的可逆性，得到
	\[
		\dot\omega_t=-\mathrm{d}(\iota_{X_t}\omega_t).
	\]
	对时间积分便有
	\[
		\omega_t-\omega_0
		=-\mathrm{d}\left(\int_0^t\iota_{X_s}\omega_s\,\mathrm{d}s\right),
	\]
	所以差值恰当，上同调类恒定。
\end{proof}

\term[Moserwendingxingdingli]{Moser 稳定性定理}[Moser stability theorem]比较的是一整条处处非退化且同上同调类的路径。给定两个端点后，直接连接的形式 \((1-t)\omega_0+t\omega_1\) 虽然总闭，却未必始终非退化；例如在四维辛空间中，\(\omega\) 与 \(-\omega\) 都非退化且给出相同定向，其直线路径仍在中点成为零。Moser 论证使用的逐点逆算子正是在退化处失去定义。整体方法及其原始形式条件可对照 Cannas da Silva 的定理1.7及证明。\cite[第1.4节，第88--90页]{Cannas2006Symplectic}

\begin{example}\label{g4:exm:0801:torus-density}
	在 \(M=(\mathbb R/\mathbb Z)^2\) 上用周期坐标 \(q,p\)，固定 \(|\varepsilon|<1\)，令
	\[
		\omega_t=(1+t\varepsilon\cos(2\pi q))\,\mathrm{d}q\wedge\mathrm{d}p,
		\qquad
		\alpha_t=\frac{\varepsilon}{2\pi}\sin(2\pi q)\,\mathrm{d}p.
	\]
	系数始终为正，且 \(\dot\omega_t=\mathrm{d}\alpha_t\)。Moser 方程给出
	\[
		X_t=-\frac{\varepsilon\sin(2\pi q)}{2\pi(1+t\varepsilon\cos(2\pi q))}\,\partial_q.
	\]
	令 \(F_t(q)=q+t\varepsilon\sin(2\pi q)/(2\pi)\)。因为 \(F_t'(q)>0\) 且 \(F_t(q+1)=F_t(q)+1\)，它诱导圆周微分同胚。对 \(F_t(q_t)=q\) 求导可知 \(q_t\) 恰满足上述速度方程，所以
	\[
		\phi_t(q,p)=(F_t^{-1}(q),p).
	\]
	直接求拉回也能验证：\(\omega_t\) 的密度因子是 \(F_t'\)，而逆映射的导数给出其倒数，两者相消，故 \(\phi_t^*\omega_t=\mathrm{d}q\wedge\mathrm{d}p\)。这里移动横坐标，恰好补偿了随位置变化的面积密度。
\end{example}

\subsection{紧性、紧支和相对保持}
\label{g4:subsec:0801:04}

上一证明中，拉回求导和收缩方程都是局部的；紧性负责把它们产生的速度场积分为全时段的整体同位。若原始形式及流的存在范围已经给定，相同计算可以用于非紧流形或子流形的邻域。

\begin{proposition}\label{g4:prop:0801:relative-flow}
	设 \(N\) 为无边界光滑流形，\(\omega_t\) 为光滑辛形式族，\(\alpha_t\) 为光滑一形式族，满足 \(\dot\omega_t=\mathrm{d}\alpha_t\)。令 \(X_t\) 由 \(\iota_{X_t}\omega_t=-\alpha_t\) 确定。
	\begin{enumerate}
		\item 若从时刻 \(0\) 出发的流 \(\phi_t\) 对开集 \(U\subset N\) 的所有初值及 \(0\le t\le1\) 都存在，则 \(\phi_t:U\to\phi_t(U)\) 是微分同胚，且 \(\phi_t^*\omega_t=\omega_0|_U\)。
		\item 若 \(S\subset N\) 是嵌入子流形，且每个 \(p\in S\) 都满足 \(\alpha_t(p)=0\in\mathrm{T}_p^*N\)，则存在含 \(S\) 的开集 \(U\)，使上述流在 \([0,1]\times U\) 上存在，并逐点固定 \(S\)。这里不要求 \(S\) 紧。
		\item 在第二项的条件下，若再有 \(\mathrm{D}X_t(p)=0\) 对所有 \(p\in S\) 成立，则 \(\mathrm{d}(\phi_t)_p\) 在整个 \(\mathrm{T}_pN\) 上为恒等映射。因为 \(X_t(p)=0\)，这里的线性化 \(\mathrm{D}X_t(p)\) 不依赖坐标选择。
	\end{enumerate}
\end{proposition}
\begin{proof}
	第一项的微分同胚结论来自演化映射的逆；拉回等式由\cref{g4:prop:0801:pullback-derivative}和收缩方程直接得到，其求导结果为零，初值为 \(\omega_0|_U\)。

	对第二项，\cref{g4:lem:0801:contraction-solution}给出 \(X_t(p)=0\)，所以每个 \(p\in S\) 的常值曲线都是定义在整个 \([0,1]\) 上的解。唯一性说明该点被所有时间映射固定。演化映射定义域的开放性说明，每个 \((t,p)\in[0,1]\times S\) 附近有一个时间与初值的乘积邻域，所有相应解在那里存在。固定 \(p\)，用有限个这样的时间邻域覆盖紧区间 \([0,1]\)，再交其初值邻域，得到 \(p\) 的开邻域 \(U_p\)，使从 \(0\) 出发的所有这些初值都能求解到每个 \(t\in[0,1]\)。令 \(U=\bigcup_{p\in S}U_p\)，即可同时包含整个 \(S\)。这个论证允许不同点采用不同大小的邻域，不要求给非紧 \(S\) 一个统一半径。

	第三项可在固定点 \(p\) 附近选坐标验证。对流方程的初值求导，令 \(A_t=\mathrm{d}(\phi_t)_p\)，由于 \(\phi_t(p)=p\)，得到
	\[
		\dot A_t=\mathrm{D}X_t(p)\,A_t,\qquad A_0=I.
	\]
	在所给条件下 \(\dot A_t=0\)，所以 \(A_t=I\)。换坐标时，向量场一阶导数的额外项带有 \(X_t(p)\) 作为因子，因此在该零点处消失，线性化的定义也就与坐标无关。
\end{proof}

相对条件必须读作整个协向量为零。若只要求包含映射 \(i:S\hookrightarrow N\) 满足 \(i^*\alpha_t=0\)，便只知道 \(\alpha_t\) 在 \(\mathrm{T}S\) 上为零，因而 \(X_t\in(\mathrm{T}S)^{\omega_t}\)；这不保证 \(X_t=0\)。例如 \(N=\mathbb R^2\)、\(\omega=\mathrm{d}q\wedge\mathrm{d}p\)、\(S=\{p=0\}\)，取常值闭一形式 \(\alpha=\mathrm{d}p\)。虽然 \(i^*\alpha=0\)，收缩方程却给出 \(X=-\partial_q\)，其流沿 \(S\) 平移，并不逐点固定 \(S\)。

\ProseParagraphBreak

在邻域问题中，通常先构造沿 \(S\) 消失的原始形式，再用上述命题缩小初值邻域。这个顺序使“固定子流形”成为微分方程的初值事实。Cannas da Silva 的相对定理1.8采用同一思路；沿子流形构造原始形式的细节将在\cref{g4:sec:0803}展开。\cite[第1.5节，第90页]{Cannas2006Symplectic}

\begin{corollary}\label{g4:cor:0801:compact-support}
	设 \(M\) 为无边界光滑流形，\(\omega_t\) 为光滑辛形式族。如果存在光滑一形式族 \(\alpha_t\) 和固定紧集 \(K\subset M\)，满足
	\[
		\dot\omega_t=\mathrm{d}\alpha_t,\qquad
		\operatorname{supp}\alpha_t\subset K\quad(0\le t\le1),
	\]
	则存在从恒等映射出发的同位 \(\phi_t:M\to M\)，满足 \(\phi_t^*\omega_t=\omega_0\)，且 \(\phi_t\) 在 \(M\setminus K\) 上恒等。
\end{corollary}
\begin{proof}
	逐点求逆给出的 \(X_t\) 在 \(K\) 外为零，故各 \(X_t\) 具有共同紧支。第一卷命题13.2.10的延拓论证保证其演化映射在整个时间区间及所有初值上存在；\(K\) 外的点由唯一性保持不动。再用\cref{g4:prop:0801:relative-flow}第一项即可。
\end{proof}

这里要求的是原始一形式具有共同紧支，而非仅仅每个固定时刻的速度场分别紧支。若没有这样的控制或其他完备性保证，非紧流形上的全局结论确实可能失败。

\begin{example}\label{g4:exm:0801:noncompact-area}
	在 \(\mathbb R^2\) 上令
	\[
		\omega_t=\bigl(1-t+t e^{-(q^2+p^2)}\bigr)\,\mathrm{d}q\wedge\mathrm{d}p,
		\qquad 0\le t\le1.
	\]
	每个系数都严格为正，所以这是一条光滑辛路径。令
	\[
		\sigma(q,p)=\left(\int_0^q(e^{-(s^2+p^2)}-1)\,\mathrm{d}s\right)\mathrm{d}p,
	\]
	则 \(\omega_t-\omega_0=t\,\mathrm{d}\sigma\)。因此整个族同上同调类，而且所需的光滑原始一形式已经显式给出。

	然而 \(\omega_0\) 的总面积无限，\(\omega_1\) 的总面积为
	\[
		\int_{\mathbb R^2}\omega_1
		=2\pi\int_0^\infty e^{-r^2}r\,\mathrm{d}r=\pi.
	\]
	若有微分同胚 \(\phi\) 满足 \(\phi^*\omega_1=\omega_0\)，则它保持这两个形式给出的定向。对半径为 \(R\) 的圆盘 \(B_R\)，换元公式给出
	\[
		\pi R^2=\int_{B_R}\omega_0
		=\int_{\phi(B_R)}\omega_1\le\pi,
	\]
	当 \(R>1\) 时矛盾。因此两个端点之间根本不存在这样的整体辛同胚。失败并非某一次原始形式选择不佳，而是非紧流形上无穷远处的面积无法被整体微分同胚压入有限总面积。
\end{example}
% 纲要条目排为编号小节标题，并进入目录与书签；说明段仍保留于纲要注释中。

% !TeX root = ../../Book4.tex
% 纲要标题：### 8.2 Darboux 定理
% 纲要位置：几何学计划纲要.md，第 3786 行。
% 正文状态：初稿；尚未正式审核。

% 纲要内容开始
% 
% * 局部线性化和直线路径
% * 非退化性的小邻域控制
% * 原始形式和径向同伦
% * 局部坐标的构造及无局部曲率不变量
% 
% 纲要内容结束

\section{Darboux 定理}
\label{g4:sec:0802}

辛基把一个切空间上的交替型化成标准形。能否找到一组实际坐标，使整个邻域中的辛形式都成为同一个常系数形式？本节用 Moser 方程完成这一步：先在一点匹配线性数据，再用径向同伦求出一个在该点二阶消失的原始一形式，最后控制由它产生的局部运动。原始形式的消失阶数还会保证，修正坐标时保留最初指定的辛标架。Darboux 定理及这种证明方法可参见 \cite[第1.4--1.5节，特别是第90--91页的定理1.8、1.9]{Cannas2006Symplectic}。

\subsection{局部线性化和直线路径}
\label{g4:subsec:0802:01}

设 \((M^{2n},\omega)\) 为光滑辛流形，\(m\in M\)。由\cref{g4:thm:0701:symplectic-basis}，可以在 \(\mathrm{T}_mM\) 选取辛基
\[
	e_1,\ldots,e_n,f_1,\ldots,f_n,\qquad
	\omega_m(e_i,f_j)=\delta_{ij}.
\]
更具体地，记
\[
	\omega_0=\sum_{i=1}^n\mathrm{d}q^i\wedge\mathrm{d}p_i
\]
为 \(\mathbb R^{2n}\) 上的标准形式，并以 \(A:\mathbb R^{2n}\to\mathrm{T}_mM\) 表示将标准坐标基送到上述辛基的线性同构。因此 \(A^*\omega_m=(\omega_0)_0\)。

\ProseParagraphBreak

这个指定的一阶数据可以由一个普通坐标参数化实现。任取 \(\kappa_0:U_0\subset\mathbb R^{2n}\to M\)，使 \(\kappa_0(0)=m\)，并令
\[
	L=(\mathrm{d}\kappa_0|_0)^{-1}A,\qquad
	\kappa=\kappa_0\circ L.
\]
缩小定义域后，\(\kappa\) 是到 \(m\) 的一个邻域的微分同胚，而且 \(\mathrm{d}\kappa|_0=A\)。把辛形式拉回到 Euclid 空间，记
\[
	\widetilde\omega=\kappa^*\omega,\qquad
	\beta=\widetilde\omega-\omega_0.
\]
则 \(\beta\) 是光滑闭二形式，并且 \(\beta_0=0\)。这里的等号指两个二形式在原点的完整交替配对相等。它暂时只匹配原点处的值；附近各点的系数仍然可能变化。

\ProseParagraphBreak

在一个以原点为中心的坐标球上取直线路径
\begin{equation}
	\label{g4:eq:0802:straight-path}
	\omega_t=\omega_0+t\beta,\qquad 0\le t\le1.
\end{equation}
每个 \(\omega_t\) 都闭，而且 \((\omega_t)_0=(\omega_0)_0\)。我们将寻找在更小球上定义、到其像的微分同胚 \(\phi_t\)，满足
\[
	\phi_0=\operatorname{id},\qquad
	\phi_t^*\omega_t=\omega_0,\qquad
	\phi_t(0)=0,\qquad \mathrm{d}\phi_t|_0=I.
\]
这样 \(\Psi=\kappa\circ\phi_1\) 就把标准形式送到原流形的辛形式，并且 \(\mathrm{d}\Psi|_0=A\)。首先要保证，式\eqref{g4:eq:0802:straight-path}的全部中间形式在同一个邻域内都非退化。

\subsection{非退化性的小邻域控制}
\label{g4:subsec:0802:02}

两个辛形式之间的直线段一般会经过退化形式。例如 \(\omega_0\) 与 \(-\omega_0\) 都非退化，而它们的中点是零。当前路径的有利条件是：原点处的全部形式已经等于同一个 \(\omega_0\)，所以可以在原点附近同时控制它们。

\begin{lemma}
	\label{g4:lem:0802:uniform-nondegeneracy}
	设 \(\beta\) 是 \(\mathbb R^{2n}\) 中原点附近的光滑二形式，\(\beta_0=0\)，并令 \(\omega_t=\omega_0+t\beta\)。则存在 \(R>0\)，使对所有 \(x\in B_R(0)\) 和 \(0\le t\le1\)，收缩映射
	\[
		b_t(x):\mathbb R^{2n}\longrightarrow(\mathbb R^{2n})^*,
		\qquad b_t(x)v=(\omega_t)_x(v,\cdot)
	\]
	可逆。采用标准 Euclid 范数时，还可使
	\[
		\|b_t(x)^{-1}\|\le2.
	\]
	逆映射依赖 \((t,x)\) 光滑，并在包含 \([0,1]\) 的开时间区间上仍有定义。
\end{lemma}
\begin{proof}
	用 Euclid 度量给向量空间及其对偶赋范数，记 \(b_0v=\omega_0(v,\cdot)\)。标准坐标中的计算给出 \(\|b_0v\|=\|v\|\)。因 \(\beta_0=0\)，可以取闭球 \(\overline B_R\) 留在原定义域内，并使
	\[
		\sup_{x\in B_R}\|\beta_x(\,\cdot\,,\,\cdot\,)\|<\frac14,
	\]
	其中二形式的范数也就是相应收缩映射的算子范数。对 \(-1\le t\le2\)，有
	\[
		\|b_t(x)v\|
		\ge\|b_0v\|-|t|\,\|\beta_x(v,\cdot)\|
		\ge\frac12\|v\|.
	\]
	因此 \(b_t(x)\) 单射；源和目标同维，所以它是同构，且逆的范数至多为 \(2\)。矩阵求逆在可逆矩阵开集上光滑，故逆映射依赖 \((t,x)\) 光滑。估计覆盖了比 \([0,1]\) 更大的时间区间，也保证了端点附近的时间延拓。
\end{proof}

这个证明同时给出了两个用途不同的结论：同一个球适用于整条路径，收缩方程的逆在这条路径上有统一界。后一个界会把原始一形式的消失阶数传给 Moser 向量场，并使我们能够估计轨道移动的距离。

\subsection{原始形式和径向同伦}
\label{g4:subsec:0802:03}

为了使用上一节的方程，需要一个满足 \(\mathrm{d}\alpha=\beta\) 的一形式。任取一个原始形式可能在原点不为零，那样得到的流就会移动原点。径向同伦利用从 \(0\) 到 \(x\) 的线段作积分，同时提供原始形式与原点处的控制。

\begin{lemma}
	\label{g4:lem:0802:radial-primitive}
	设 \(\beta\) 是球 \(B_R(0)\subset\mathbb R^{2n}\) 上的光滑闭二形式，且 \(\beta_0=0\)。定义
	\begin{equation}
		\label{g4:eq:0802:radial-primitive}
		\alpha_x(v)=\int_0^1 s\,\beta_{sx}(x,v)\,\mathrm{d}s,
		\qquad x\in B_R,\quad v\in\mathbb R^{2n}.
	\end{equation}
	则 \(\alpha\) 光滑，满足 \(\mathrm{d}\alpha=\beta\)。缩小 \(R\) 后，存在常数 \(C_0\)，使
	\[
		\|\alpha_x\|\le C_0|x|^2.
	\]
	特别地，\(\alpha\) 的全部坐标系数及其一阶导数在原点都为零。
\end{lemma}
\begin{proof}
	式\eqref{g4:eq:0802:radial-primitive}是第一卷定理14.2.2中径向同伦算子在二形式上的公式。其因子 \(s\) 来自一个切向量被缩放映射 \(x\mapsto sx\) 送成 \(sv\)，另一个输入 \(x\) 则是径向运动的速度。被积式在 \([0,1]\times B_R\) 上光滑，因此参数积分给出光滑一形式。

	为核对 \(\mathrm{d}\alpha=\beta\) 的符号，合写坐标为 \(x^1,\ldots,x^{2n}\)，并置
	\[
		\beta=\frac12\sum_{i,j}b_{ij}(x)\,\mathrm{d}x^i\wedge\mathrm{d}x^j,
		\qquad b_{ij}=-b_{ji}.
	\]
	式\eqref{g4:eq:0802:radial-primitive}给出
	\[
		\alpha=\sum_j a_j(x)\,\mathrm{d}x^j,\qquad
		a_j(x)=\sum_i x^i\int_0^1s\,b_{ij}(sx)\,\mathrm{d}s.
	\]
	闭性 \(\mathrm{d}\beta=0\) 在分量上等价于
	\(\partial_kb_{ij}-\partial_jb_{ik}=\partial_i b_{kj}\)。于是
	\[
		\begin{aligned}
			\partial_k a_j-\partial_j a_k
			&=\int_0^1\left(2s\,b_{kj}(sx)
				+s^2\sum_i x^i\partial_i b_{kj}(sx)\right)\mathrm{d}s\\
			&=\int_0^1\frac{\mathrm{d}}{\mathrm{d}s}
				\left(s^2b_{kj}(sx)\right)\mathrm{d}s=b_{kj}(x).
		\end{aligned}
	\]
	这证明了 \(\mathrm{d}\alpha=\beta\)。

	最后，\(\beta_0=0\) 且 \(\beta\) 光滑，故在一个较小闭球上可取常数 \(C\) 使 \(\|\beta_y\|\le C|y|\)。从径向公式直接得到
	\[
		|\alpha_x(v)|
		\le\int_0^1s\,C|sx|\,|x|\,|v|\,\mathrm{d}s
		=\frac C3|x|^2|v|.
	\]
	令 \(C_0=C/3\) 即得估计。因为 \(\alpha\) 光滑且其系数为 \(O(|x|^2)\)，这些系数及其一阶导数在原点全部为零。
\end{proof}

现在在\cref{g4:lem:0802:uniform-nondegeneracy}的共同球上，按 Moser 方程定义
\begin{equation}
	\label{g4:eq:0802:moser-field}
	\iota_{X_t}\omega_t=-\alpha,
	\qquad
	X_t(x)=-b_t(x)^{-1}\alpha_x.
\end{equation}
光滑性与唯一性来自收缩映射的可逆性，逆的统一界则给出
\begin{equation}
	\label{g4:eq:0802:field-estimate}
	|X_t(x)|\le C_1|x|^2,\qquad
	X_t(0)=0,\qquad
	\mathrm{D}_xX_t(0)=0,
	\qquad 0\le t\le1.
\end{equation}
其中可取 \(C_1=2C_0\)。最后一式由光滑乘积 \(b_t^{-1}\alpha\) 的求导及 \(\alpha\) 的一阶消失得到。这比只知道 \(\mathrm{d}\alpha|_0=0\) 更强：外微分为零只消去导数的交替部分，而这里全部一阶系数都消失。

\subsection{局部坐标的构造及无局部曲率不变量}
\label{g4:subsec:0802:04}

\term[Darbouxdingli]{Darboux 定理}[Darboux theorem]\footnote{Jean Gaston Darboux（1842--1917），法国数学家，研究微分几何与分析。}把上述构造汇集为局部坐标的存在性，并允许预先指定原点处的辛标架。

\begin{theorem}[Darboux 定理]
	\label{g4:thm:0802:darboux}
	设 \((M^{2n},\omega)\) 为光滑无边界辛流形，\(m\in M\)，并给定辛线性同构
	\[
		A:(\mathbb R^{2n},\omega_0)\longrightarrow
		(\mathrm{T}_mM,\omega_m).
	\]
	则存在 \(0\) 的开邻域 \(V\subset\mathbb R^{2n}\)、\(m\) 的开邻域 \(U\subset M\)，以及微分同胚 \(\Psi:V\to U\)，满足
	\[
		\Psi(0)=m,\qquad \mathrm{d}\Psi|_0=A,\qquad
		\Psi^*\omega=\omega_0.
	\]
	等价地，\(U\) 上存在以 \(m\) 为原点的坐标 \((q^1,\ldots,q^n,p_1,\ldots,p_n)\)，使
	\[
		\omega=\sum_{i=1}^n\mathrm{d}q^i\wedge\mathrm{d}p_i,
	\]
	且原点处的坐标切向量实现指定的辛标架。这些坐标称为\term[Darbouxzuobiao]{Darboux 坐标}[Darboux coordinates]。
\end{theorem}
\begin{proof}
	按第一个小节选取 \(\kappa\)，使 \(\kappa(0)=m\)、\(\mathrm{d}\kappa|_0=A\)，并在坐标球上令 \(\widetilde\omega=\kappa^*\omega\)、\(\beta=\widetilde\omega-\omega_0\)。前两个引理给出 \(\omega_t=\omega_0+t\beta\) 的共同非退化球 \(B_R\)，以及 \(\mathrm{d}\alpha=\beta\)、\(\alpha=O(|x|^2)\)。由式\eqref{g4:eq:0802:moser-field}定义 \(X_t\)，它满足式\eqref{g4:eq:0802:field-estimate}。

	先控制流的存在时间。第一卷命题13.1.9给出光滑时间依赖场的局部演化映射。由于 \(X_t\) 只在 \(B_R\) 中定义，我们在其中选两个更小的球：取 \(r>0\)，使
	\[
		2r<R,\qquad 4C_1r<\frac12.
	\]
	若 \(C_1=0\)，第二个条件自动满足。令 \(x(t)\) 从 \(x_0\in B_r\) 出发，满足 \(\dot x=X_t(x)\)。只要 \(0\le t\le1\) 且轨道尚在 \(B_{2r}\) 中，就有
	\[
		|x(t)-x_0|
		\le\int_0^t C_1|x(s)|^2\,\mathrm{d}s
		\le4C_1r^2t<\frac r2,
	\]
	从而 \(|x(t)|<3r/2\)。因此轨道不能在 \(t\le1\) 首次到达 \(\partial B_{2r}\)。

	还要排除在到达时间 \(1\) 之前终止的可能。如果最大解有这样的有限端点，将时间并入状态，考虑增维向量场 \(\partial_t+X_t\)。上述估计使相应轨道留在固定紧集 \([0,1]\times\overline B_{3r/2}\) 中；这个紧集位于增维场的开定义域内，因为 \(X_t\) 在 \([0,1]\) 两端附近仍光滑。第一卷定理13.2.2的有限端点逃逸判据便允许继续延拓，与最大性矛盾。所以每个 \(x_0\in B_r\) 的解均在全部 \(0\le t\le1\) 上存在。记其演化为 \(\phi_t(x_0)\)。第一卷命题13.1.9还说明，\(\phi_t\) 光滑，并由反向演化给出到其开像的光滑逆。

	因为 \(X_t(0)=0\)，唯一性给出 \(\phi_t(0)=0\)。第一卷定理13.1.6及其变分方程说明 \(B_t=\mathrm{d}\phi_t|_0\) 满足
	\[
		\dot B_t=\mathrm{D}_xX_t(\phi_t(0))B_t=0,
		\qquad B_0=I.
	\]
	因此 \(\mathrm{d}\phi_t|_0=I\)，初始辛标架得到保留。

	现在应用\cref{g4:prop:0801:pullback-derivative}。它只要求当前演化在所考察的开集和时间区间上存在，这一条件已经验证。由于 \(\dot\omega_t=\beta=\mathrm{d}\alpha\)、\(\mathrm{d}\omega_t=0\)，有
	\[
		\begin{aligned}
			\frac{\mathrm{d}}{\mathrm{d}t}(\phi_t^*\omega_t)
			&=\phi_t^*\bigl(\dot\omega_t+\mathcal L_{X_t}\omega_t\bigr)\\
			&=\phi_t^*\mathrm{d}\bigl(\alpha+\iota_{X_t}\omega_t\bigr)=0.
		\end{aligned}
	\]
	初值为 \(\omega_0\)，所以 \(\phi_1^*\widetilde\omega=\omega_0\)。取 \(V=B_r\)、\(U=\kappa(\phi_1(B_r))\)、\(\Psi=\kappa\circ\phi_1\)，便得到
	\[
		\Psi^*\omega=\phi_1^*(\kappa^*\omega)=\omega_0,\qquad
		\mathrm{d}\Psi|_0=\mathrm{d}\kappa|_0\,\mathrm{d}\phi_1|_0=A.
	\]
	逆映射 \(\Psi^{-1}\) 的分量就是所需坐标。
\end{proof}

% !TeX root = ../Book4.tex
\begin{figure}[htbp]
	\centering
	\begin{tikzpicture}[font=\small,scale=0.9]
		\fill[GeometryGrayLight] (0,0) circle (2.6);
		\fill[GeometryBlueLight] (0,0) circle (1);
		\draw[GeometryGray,thick] (0,0) circle (2.6);
		\draw[GeometryGray,dashed] (0,0) circle (2);
		\draw[GeometryBlue] (0,0) circle (1);
		\fill (0,0) circle (1.6pt);
		\node[below left] at (0,0) {\(0\)};
		\node[below] at (0,-1) {\(B_r\)};
		\node[below] at (0,-2) {\(B_{2r}\)};
		\node[below] at (0,-2.6) {\(B_R\)};
		\fill[GeometryBlue] (0.65,0.35) circle (2pt);
		\node[left] at (0.6,0.4) {\(x_0\)};
		\draw[GeometryBlue,thick,->] (0.65,0.35)
			.. controls (0.75,0.5) and (0.85,0.58) .. (1,0.55);
		\node[above right] at (1,0.55) {\(\phi_t(x_0)\)};
	\end{tikzpicture}
	\caption[Darboux 构造中的局部流范围]{在 \(B_R\) 上构造向量场，从更小的 \(B_r\) 出发。二次速度估计保证轨道在整个单位时间内留在 \(B_{2r}\) 中。图为嵌套 Euclid 球及一条轨道的二维示意，曲线形状不代表特定方程的解。}
	\label{g4:fig:0802:flow-domains}
\end{figure}


\cref{g4:fig:0802:flow-domains}表示证明中三层邻域的作用。原始形式和向量场在 \(B_R\) 上构造，实际坐标从 \(B_r\) 中的初值产生，中间的 \(B_{2r}\) 则用来验证轨道不会在时间 \(1\) 以前离开构造区域。这个大小控制把局部常微分方程接到了 Moser 拉回恒等式上。

\begin{example}
	\label{g4:exm:0802:surface-coordinates}
	在平面点 \((x_0,y_0)\) 附近，设
	\[
		\omega=f(x,y)\,\mathrm{d}x\wedge\mathrm{d}y,\qquad
		f(x_0,y_0)\ne0.
	\]
	缩小为一个 \(f\) 不为零的矩形，取
	\[
		Q=x-x_0,\qquad
		P(x,y)=\int_{y_0}^{y}f(x,s)\,\mathrm{d}s.
	\]
	则
	\[
		\mathrm{d}Q\wedge\mathrm{d}P
		=\mathrm{d}x\wedge
		\left(\partial_xP\,\mathrm{d}x+f(x,y)\,\mathrm{d}y\right)
		=f(x,y)\,\mathrm{d}x\wedge\mathrm{d}y.
	\]
	坐标变换的 Jacobian 行列式为 \(f(x,y)\ne0\)，所以第一卷定理11.3.3的逆函数定理给出以该点为中心的 Darboux 坐标。例如 \(f=1+x^2+y^2\)、\((x_0,y_0)=(0,0)\) 时，
	\[
		Q=x,\qquad P=(1+x^2)y+\frac{y^3}{3}.
	\]
	非恒定的面积系数被吸收到新坐标 \(P\) 中。
\end{example}

\begin{example}
	\label{g4:exm:0802:separated-four-dimensional}
	在 \(\mathbb R^4\) 的坐标 \((x_1,y_1,x_2,y_2)\) 中，考虑
	\[
		\omega=(1+x_1^2+y_1^2)\,\mathrm{d}x_1\wedge\mathrm{d}y_1
		+e^{x_2}\,\mathrm{d}x_2\wedge\mathrm{d}y_2.
	\]
	每项的系数只依赖本项涉及的两个坐标，所以求外微分时产生重复的一形式，\(\mathrm{d}\omega=0\)。两个系数都为正，因此逐对配对也显示其非退化。分别取
	\[
		Q_1=x_1,\quad P_1=(1+x_1^2)y_1+\frac{y_1^3}{3},
		\qquad
		Q_2=x_2,\quad P_2=e^{x_2}y_2.
	\]
	直接求导即得
	\[
		\omega=\mathrm{d}Q_1\wedge\mathrm{d}P_1
		+\mathrm{d}Q_2\wedge\mathrm{d}P_2.
	\]
	按配对次序排列坐标，变换的 Jacobian 行列式为
	\((1+x_1^2+y_1^2)e^{x_2}>0\)，故这些函数在每一点附近都是坐标。这个可分计算在四维中同时整理了两对方向；一般辛形式具有混合项时，定理的局部流构造统一完成同样的任务。
\end{example}

Darboux 定理说明，同一维数的辛形式在每一点附近都与同一个标准形式等价。若一个点态量只由辛形式的局部数据决定，并且在局部辛同胚下保持不变，那么它在所有同维辛流形的各点都必须取相同的值。因此，辛形式本身没有像黎曼曲率那样随点变化、阻碍局部标准化的曲率不变量。这里的结论允许任意局部坐标变化；具体坐标中出现变化的系数，并不因此成为几何不变量。

\ProseParagraphBreak

附加结构仍然可以保留局部信息。若另给一个黎曼度量，度量的曲率属于额外数据，Darboux 坐标并不同时使它成为常系数。若另给一个相容几乎复结构 \(J\)，可积性也要另外检验：\cref{g4:exm:0704:nonintegrable-compatible}中的辛形式已经是常系数形式，但所给 \(J\) 处处不可积。局部标准形整理的是辛形式；同时保持度量或几乎复结构，会提出更强的等价要求。
% 纲要条目排为编号小节标题，并进入目录与书签；说明段仍保留于纲要注释中。

% !TeX root = ../../Book4.tex
% 纲要标题：### 8.3 沿子流形的相对方法
% 纲要位置：几何学计划纲要.md，第 3793 行。
% 正文状态：初稿；尚未正式审核。

% 纲要内容开始
% 
% * 相对 Poincaré 引理
% * 匹配法向线性数据
% * 辛子流形邻域定理
% * Lagrange 邻域定理
% 
% 纲要内容结束

\section{沿子流形的相对方法}
\label{g4:sec:0803}

Darboux 坐标把一个点附近的辛形式化为标准形式。若要同时保留一整条子流形，坐标改变便受到约束：不仅要把形式变对，还要知道子流形上的点及横向的一阶数据如何变化。本节把这两项任务分开。先用管状坐标匹配沿子流形的线性数据，再用一个沿子流形消失的向量场修正邻域中的形式。辛子流形由自身的辛形式与辛法丛控制邻域；Lagrange 子流形则以余切丛为统一模型。

\subsection{相对 Poincaré 引理}
\label{g4:subsec:0803:01}

设 \(i:S\hookrightarrow M\) 为嵌入。对二形式 \(\beta\)，条件 \(i^*\beta=0\) 只表示 \(\beta_x(u,v)=0\) 对 \(u,v\in\mathrm{T}_xS\) 成立；更强的条件 \(\beta_x=0\) 表示它在整个 \(\mathrm{T}_xM\times\mathrm{T}_xM\) 上为零。前者使邻域中的闭形式有相对原始，后者还将使原始的一阶导数消失。下面逐项证明这两个结论。

\ProseParagraphBreak

第二卷定理15.3.5给出子流形的\term[guanzhuanglinyu]{管状邻域}[tubular neighborhood]：可将 \(S\) 的一个邻域识别为某向量丛 \(\pi:E\to S\) 的变半径圆盘丛
\[
	W=\{e\in E:|e|<\varepsilon(\pi(e))\},\qquad \varepsilon(x)>0.
\]
零截面记为 \(s:S\to W\)。这样的邻域在纤维缩放 \(H_t(e)=te\)、\(0\le t\le1\) 下保持不变；\(H_1=\operatorname{id}\)，\(H_0=s\circ\pi\)。半径随基点变化不影响这一性质。

\begin{lemma}[相对 Poincaré 引理]\label{g4:lem:0803:relative-poincare}
	设 \(\pi:E\to S\) 是无边界光滑流形上的实光滑向量丛，\(W\subset E\) 是包含零截面、在 \(e\mapsto te\)、\(0\le t\le1\) 下不变的开集。对 \(k\ge1\) 及 \(\gamma\in\Omega^k(W)\)，令
	\begin{equation}\label{g4:eq:0803:fiber-homotopy}
		\begin{aligned}
			&(K\gamma)_e(v_1,\ldots,v_{k-1})\\
			&\quad=\int_0^1\gamma_{te}\bigl(\partial_tH(t,e),\mathrm{d}H_t(v_1),\ldots,\mathrm{d}H_t(v_{k-1})\bigr)\,\mathrm{d}t.
		\end{aligned}
	\end{equation}
	其中 \(H(t,e)=te\)。则 \(K\gamma\) 光滑，并满足
	\begin{equation}\label{g4:eq:0803:homotopy-identity}
		\mathrm{d}K\gamma+K\mathrm{d}\gamma=\gamma-\pi^*s^*\gamma.
	\end{equation}
	特别地，若 \(\beta\in\Omega^2(W)\) 闭且 \(s^*\beta=0\)，则 \(\alpha=K\beta\) 满足 \(\mathrm{d}\alpha=\beta\)，且 \(\alpha_{0_x}=0\) 作为整个 \(\mathrm{T}_{0_x}^*E\) 中的协向量成立。若进一步有 \(\beta_{0_x}=0\) 作为整个环境切空间上的二形式成立，则 \(\alpha\) 的所有局部系数及其全部一阶偏导数沿零截面都为零。
\end{lemma}
\begin{proof}
	映射 \(H:[0,1]\times W\to W\) 光滑，式\eqref{g4:eq:0803:fiber-homotopy}的被积函数在 \(t=0\) 也光滑，故对紧参数区间积分仍得到光滑形式。将拉回分解为
	\[
		H^*\gamma=\mathrm{d}t\wedge\eta_t+\zeta_t,
		\qquad \zeta_t=H_t^*\gamma,
	\]
	其中 \(\eta_t\) 与 \(\zeta_t\) 都只含 \(W\) 方向的微分，且 \(K\gamma=\int_0^1\eta_t\,\mathrm{d}t\)。外微分给出
	\[
		H^*\mathrm{d}\gamma
		=\mathrm{d}t\wedge\bigl(\partial_t\zeta_t-\mathrm{d}\eta_t\bigr)+\mathrm{d}\zeta_t.
	\]
	因此
	\[
		K\mathrm{d}\gamma
		=\zeta_1-\zeta_0-\mathrm{d}\int_0^1\eta_t\,\mathrm{d}t
		=\gamma-\pi^*s^*\gamma-\mathrm{d}K\gamma,
	\]
	证明同伦恒等式。若 \(\mathrm{d}\beta=0\) 且 \(s^*\beta=0\)，立即得到 \(\mathrm{d}K\beta=\beta\)。

	在零向量 \(e=0_x\) 处，路径 \(H(t,0_x)\) 恒定，故 \(\partial_tH(t,0_x)=0\)。式\eqref{g4:eq:0803:fiber-homotopy}于是对任意环境切向量都取零，证明 \(\alpha_{0_x}=0\)。这比仅有 \(s^*\alpha=0\) 强，因为后者只检验切于零截面的向量。

	最后取局部丛坐标 \((x,y)\)，其中 \(y=0\) 为零截面。若 \(\beta_{(x,0)}=0\)，则它相对于 \(\mathrm{d}x^i,\mathrm{d}y^a\) 的所有系数 \(b(x,y)\) 均满足 \(b(x,0)=0\)，并可写成
	\[
		b(x,y)=\sum_a y^a\int_0^1\partial_{y^a}b(x,sy)\,\mathrm{d}s.
	\]
	在 \(K\beta\) 的公式中，\(\partial_tH=(0,y)\) 再贡献一个纤维坐标因子，而 \(\mathrm{d}H_t(u,v)=(u,tv)\)。因此 \(\alpha\) 的每个系数都是若干 \(y^ay^b\) 乘光滑函数的和。其值及所有一阶偏导数在 \(y=0\) 时均为零。由于系数本身已为零，这一阶消失性质在换坐标后仍成立。
\end{proof}

\begin{example}\label{g4:exm:0803:relative-vanishing}
	在 \(\mathbb R^2_{x,y}\) 中取 \(S=\{y=0\}\)、\(\beta=\mathrm{d}x\wedge\mathrm{d}y\)，并沿 \(y\) 方向缩放。虽然 \(i^*\beta=0\)，\(\beta\) 沿 \(S\) 作为环境二形式并不为零。公式给出 \(K\beta=-y\,\mathrm{d}x\)：它作为整个协向量沿 \(S\) 为零，但一阶导数并不为零，且 \(\mathrm{d}(K\beta)=\beta\)。这一算例分别展示了引理中两种消失条件的作用。
\end{example}

把引理通过管状邻域映射送回 \(M\)，便得到沿闭嵌入子流形的相对原始。对 Moser 方程，原始的消失使向量场固定子流形；原始的一阶消失还使流的微分在那里保持恒等。

\begin{theorem}[相对 Moser 定理]\label{g4:thm:0803:relative-moser}
	设 \(M\) 为无边界光滑流形，\(S\subset M\) 为闭的无边界嵌入子流形，\(\omega_0,\omega_1\) 为 \(S\) 的某共同开邻域上的辛形式。假设对每个 \(x\in S\)，有
	\begin{equation}\label{g4:eq:0803:ambient-agreement}
		(\omega_0)_x(u,v)=(\omega_1)_x(u,v)
		\qquad\bigl(u,v\in\mathrm{T}_xM\bigr).
	\end{equation}
	则存在含 \(S\) 的开邻域 \(U\) 及光滑族 \(\varphi_t:U\to M\)、\(0\le t\le1\)，每个 \(\varphi_t\) 都是到其开像的微分同胚，且
	\[
		\varphi_0=\operatorname{id},\qquad
		\varphi_t|_S=\operatorname{id}_S,\qquad
		\mathrm{d}\varphi_t|_{\mathrm{T}M|_S}=\operatorname{id},\qquad
		\varphi_t^*((1-t)\omega_0+t\omega_1)=\omega_0.
	\]
	特别地，\(\varphi_1\) 是固定 \(S\) 且保持上述一阶数据的邻域辛同胚。不要求 \(S\) 紧致。
\end{theorem}
\begin{proof}
	令 \(\beta=\omega_1-\omega_0\)、\(\omega_t=\omega_0+t\beta\)。两形式闭，故 \(\beta\) 闭；式\eqref{g4:eq:0803:ambient-agreement}说明 \(\beta\) 沿 \(S\) 作为整个二张量为零。在 \([0,1]\times S\) 上，\(\omega_t\) 始终等于非退化的 \(\omega_0\)。对每个 \(x\in S\)，用非退化性的开性及参数区间的有限覆盖，可选一个点邻域 \(V_x\)，使所有 \(\omega_t\) 在 \(V_x\) 上都非退化。取这些邻域的并，再在其中选一张变半径管状邻域，便得到所有 \(t\) 共用的非退化区间与空间定义域 \(V\)。

	由\cref{g4:lem:0803:relative-poincare}，在缩小后的 \(V\) 上有光滑一形式 \(\alpha\)，满足 \(\mathrm{d}\alpha=\beta\)，且 \(\alpha\) 的值与一阶导数沿 \(S\) 为零。逐点求解
	\[
		\iota_{X_t}\omega_t=-\alpha.
	\]
	非退化性给出唯一光滑的时间依赖向量场。由 \(\alpha|_S=0\) 得 \(X_t|_S=0\)。在坐标中对这个线性方程求一阶导数，含 \(X_t\) 的系数导数项在 \(S\) 上为零，右边的导数也为零；再次由非退化性，得到 \(\mathrm{D}X_t|_S=0\)。向量场在该点为零，所以这个导数不依赖坐标选择。

	现在处理流的存在时间。对每个 \(x\in S\)，常值曲线 \(t\mapsto x\) 是存在于整个 \([0,1]\) 的解。由\cref{g4:prop:0801:relative-flow}中的演化域开性，存在 \(x\) 的一个开邻域 \(U_x\)，其初值解都能在 \(V\) 内存在到时间1。具体地，可沿常值轨道对紧时间区间作有限次局部存在与初值连续依赖的接续。取 \(U=\bigcup_{x\in S}U_x\)，这些局部解由唯一性相容，给出同一光滑流 \(\varphi_t\)。每张映射的逆由反向演化给出，故是到开像的微分同胚。非紧时这里只逐点选择 \(U_x\)，没有使用统一正半径。

	\cref{g4:prop:0801:pullback-derivative}及闭性给出
	\[
		\frac{\mathrm{d}}{\mathrm{d}t}\varphi_t^*\omega_t
		=\varphi_t^*\bigl(\mathrm{d}\alpha+\mathrm{d}\iota_{X_t}\omega_t\bigr)=0.
	\]
	初值为 \(\omega_0\)，故所需拉回恒等式成立。常值解给出 \(\varphi_t|_S=\operatorname{id}_S\)；沿这些解的变分方程为
	\[
		\frac{\mathrm{d}}{\mathrm{d}t}\mathrm{d}\varphi_t|_x
		=\mathrm{D}X_t|_x\circ\mathrm{d}\varphi_t|_x=0,
	\]
	结合初值恒等，得到整个环境切空间上的微分恒等。
\end{proof}

若只知道 \(i^*\omega_0=i^*\omega_1\)，直线路径未必可用。例如在平面中沿 \(x\) 轴，两形式 \(\mathrm{d}x\wedge\mathrm{d}y\) 与 \(-\mathrm{d}x\wedge\mathrm{d}y\) 的拉回都为零，而其中点为零形式。接下来先匹配横向线性数据，正是为了把这种较弱的信息变成式\eqref{g4:eq:0803:ambient-agreement}。相对 Moser 的紧子流形版本见\cite[定理1.8，第90页]{Cannas2006Symplectic}；上面的变半径与流的定义域论证同时处理闭的非紧子流形。

\subsection{匹配法向线性数据}
\label{g4:subsec:0803:02}

两个子流形的邻域可能位于不同流形中。为了比较二形式，先要有一张普通微分同胚把它们放到同一空间里。只规定子流形上的映射还不够；我们希望同时实现一份预先指定的环境切丛同构。

\begin{lemma}\label{g4:lem:0803:realize-first-jet}
	设 \(M_0,M_1\) 是同维的无边界光滑流形，\(S_i\subset M_i\) 是闭的无边界嵌入子流形，\(f:S_0\to S_1\) 为微分同胚。设有覆盖 \(f\) 的光滑实向量丛同构
	\[
		F:\mathrm{T}M_0|_{S_0}\longrightarrow\mathrm{T}M_1|_{S_1},
		\qquad F|_{\mathrm{T}S_0}=\mathrm{d}f.
	\]
	则存在邻域微分同胚 \(\Phi:U_0\to U_1\)，其中 \(S_i\subset U_i\subset M_i\)，满足 \(\Phi|_{S_0}=f\) 及 \(\mathrm{d}\Phi|_{\mathrm{T}M_0|_{S_0}}=F\)。
\end{lemma}
\begin{proof}
	由第一卷定理12.3.11取一个光滑丛度量，再由该卷命题12.3.12在 \(\mathrm{T}M_0|_{S_0}\) 中取 \(\mathrm{T}S_0\) 的光滑正交补丛 \(E_0\)，并令 \(E_1=F(E_0)\)。由 \(F|_{\mathrm{T}S_0}=\mathrm{d}f\)，有
	\[
		\mathrm{T}M_i|_{S_i}=\mathrm{T}S_i\oplus E_i,
	\]
	且 \(F_E=F|_{E_0}:E_0\to E_1\) 是覆盖 \(f\) 的丛同构。

	第二卷命题15.3.1在证明管状邻域定理时计算了局部指数映射的零截面微分：若 \(\tau_i(e)=\exp_{\pi(e)}(e)\)，其中使用任意光滑联络，则
	\[
		\tau_i(0_x)=x,\qquad
		\mathrm{d}(\tau_i)_{0_x}(u,e)=u+e
		\quad\bigl(u\in\mathrm{T}_xS_i,\ e\in(E_i)_x\bigr).
	\]
	该计算只用初速度的一阶变分，因而也适用于这里选定的补丛，并不要求它是某个预定度量的法丛。右边为同构。第二卷引理15.3.3的零截面缩小论证遂给出变半径邻域 \(W_i\subset E_i\)，使 \(\tau_i:W_i\to M_i\) 都是到开像的微分同胚。

	将 \(W_0\) 换成 \(W_0\cap F_E^{-1}(W_1)\)，定义
	\begin{equation}\label{g4:eq:0803:tubular-matching}
		\Phi=\tau_1\circ F_E\circ\tau_0^{-1}.
	\end{equation}
	下面的交换图显示了这一步：上方只匹配线性纤维，下方得到实际邻域映射。
	\[
		\begin{tikzcd}[column sep=large,row sep=large]
			W_0 \arrow[r,"F_E"] \arrow[d,"\tau_0"'] & F_E(W_0) \arrow[d,"\tau_1"] \\
			U_0 \arrow[r,"\Phi"'] & U_1
		\end{tikzcd}
	\]
	零截面处 \(\mathrm{d}F_E(u,e)=(\mathrm{d}f(u),F_E(e))\)。结合两张 \(\tau_i\) 的微分公式，式\eqref{g4:eq:0803:tubular-matching}的微分把 \(u+e\) 送到 \(\mathrm{d}f(u)+F_E(e)=F(u+e)\)。因此基点映射和全部一阶数据都得到实现。
\end{proof}

这一构造也解释了法丛为何出现：选择补丛，是给商丛 \(\mathrm{T}M|_S/\mathrm{T}S\) 的每个类选择一个光滑代表；指数映射再把短代表向量变成实际的邻域点。不同选择改变初始的 \(\Phi\)，但只要它的一阶数据使辛形式匹配，相对方法就能完成修正。

\begin{corollary}\label{g4:cor:0803:symplectic-first-jet}
	设 \((M_i,\omega_i)\) 是同维无边界辛流形，\(S_i\subset M_i\) 是闭的无边界嵌入子流形。若微分同胚 \(f:S_0\to S_1\) 有覆盖它的光滑丛同构 \(F:\mathrm{T}M_0|_{S_0}\to\mathrm{T}M_1|_{S_1}\)，满足
	\[
		F|_{\mathrm{T}S_0}=\mathrm{d}f,
		\qquad (\omega_1)_{f(x)}(F_xu,F_xv)=(\omega_0)_x(u,v),
	\]
	则有邻域辛同胚 \(\Psi:U_0\to U_1\)，使 \(\Psi|_{S_0}=f\) 且 \(\mathrm{d}\Psi|_{\mathrm{T}M_0|_{S_0}}=F\)。
\end{corollary}
\begin{proof}
	由\cref{g4:lem:0803:realize-first-jet}取得 \(\Phi\)。两辛形式 \(\omega_0\) 与 \(\widetilde\omega_1=\Phi^*\omega_1\) 在 \(S_0\) 上作为整个环境二张量相等，因为 \(\mathrm{d}\Phi=F\)。\cref{g4:thm:0803:relative-moser}给出缩小邻域后的 \(\varphi\)，满足 \(\varphi^*\widetilde\omega_1=\omega_0\)、\(\varphi|_{S_0}=\operatorname{id}\)、\(\mathrm{d}\varphi|_{S_0}=I\)。取 \(\Psi=\Phi\circ\varphi\)，则
	\[
		\Psi^*\omega_1=\omega_0,\qquad
		\Psi|_{S_0}=f,\qquad
		\mathrm{d}\Psi|_{S_0}=F\circ I=F.
	\]
	令目标邻域为 \(\Psi(U_0)\) 即可。
\end{proof}

\subsection{辛子流形邻域定理}
\label{g4:subsec:0803:03}

若 \(S\subset(M,\omega)\) 本身是辛子流形，切空间的辛正交补就是一个自然的横向选择。它不需要辅助度量，而且仍带有非退化辛配对。

\begin{definition}\label{g4:def:0803:symplectic-normal}
	带有光滑纤维非退化交替型 \(\Omega\) 的实光滑向量丛 \((E,\Omega)\)，称为\term[xinxiangliangcong]{辛向量丛}[symplectic vector bundle]。辛向量丛同构是逐纤维保持这些配对的光滑实丛同构。

	设 \(S\) 为辛流形 \((M,\omega)\) 的辛子流形。子丛
	\[
		\nu_\omega S=(\mathrm{T}S)^\omega
		=\{v\in\mathrm{T}M|_S:\omega(v,u)=0\text{ 对所有 }u\in\mathrm{T}S\}
	\]
	连同限制配对 \(\omega|_{\nu_\omega S}\)，称为 \(S\) 的\term[xinfacong]{辛法丛}[symplectic normal bundle]。
\end{definition}

为核对定义中的“丛”和“辛”，考察光滑丛映射 \(v\mapsto\omega(v,\cdot)|_{\mathrm{T}S}\)。它在 \(\mathrm{T}S\) 上已为同构，因而核具有常秩，并给出光滑分裂
\begin{equation}\label{g4:eq:0803:symplectic-normal-splitting}
	\mathrm{T}M|_S=\mathrm{T}S\oplus\nu_\omega S.
\end{equation}
若 \(v\in\nu_\omega S\) 与该补空间也处处配对为零，它便与整个 \(\mathrm{T}M|_S\) 配对为零，故 \(v=0\)。因此限制形式非退化。商映射把 \(\nu_\omega S\) 同构地送到通常的法丛 \(\mathrm{T}M|_S/\mathrm{T}S\)，并在这个商上给出辛向量丛结构。

\begin{theorem}[辛子流形邻域定理]\label{g4:thm:0803:symplectic-neighborhood}
	设 \((M_i,\omega_i)\) 为无边界辛流形，\(S_i\subset M_i\) 为闭的无边界辛子流形，\(i=0,1\)。假设有辛同胚
	\[
		f:(S_0,\omega_0|_{S_0})\longrightarrow(S_1,\omega_1|_{S_1})
	\]
	以及覆盖 \(f\) 的光滑辛向量丛同构 \(A:\nu_{\omega_0}S_0\to\nu_{\omega_1}S_1\)。则存在邻域辛同胚 \(\Psi:U_0\to U_1\)，满足
	\[
		\Psi|_{S_0}=f,\qquad
		\mathrm{d}\Psi|_{\nu_{\omega_0}S_0}=A.
	\]
	这里 \(S_i\subset U_i\subset M_i\)，且子流形不必紧致。
\end{theorem}
\begin{proof}
	在式\eqref{g4:eq:0803:symplectic-normal-splitting}的两组分裂下定义
	\[
		F=\mathrm{d}f\oplus A.
	\]
	\(f\) 与 \(A\) 都保持各自的辛配对，而切向与辛法向之间的混合配对为零。因此对 \(u,v\in\mathrm{T}_xS_0\)、\(a,b\in(\nu_{\omega_0}S_0)_x\)，有
	\begin{align*}
		\omega_1\bigl(F(u+a),F(v+b)\bigr)
		&=\omega_1(\mathrm{d}f(u),\mathrm{d}f(v))+\omega_1(Aa,Ab)\\
		&=\omega_0(u,v)+\omega_0(a,b)
		=\omega_0(u+a,v+b).
	\end{align*}
	所以 \(F\) 在整个环境切丛上为辛同构，并限制为 \(\mathrm{d}f\)。\cref{g4:cor:0803:symplectic-first-jet}给出 \(\Psi\)，而其微分等于 \(F\)，特别在辛法丛上等于指定的 \(A\)。
\end{proof}

因此，辛子流形的邻域除了自身的辛几何，还记录辛法丛的拼接；若这两份数据匹配，便没有额外的邻域差别。定理的紧版本及这种法丛表述可对照\cite[定理1.11，第91--92页]{Cannas2006Symplectic}。

\begin{example}\label{g4:exm:0803:product-neighborhood}
	设 \((S,\sigma)\) 为无边界辛流形，\((N,\eta)\) 为 \(2r\) 维无边界辛流形，\(p\in N\)。乘积中的 \(S\times\{p\}\) 是辛子流形，其辛法丛为平凡丛 \(S\times\mathrm{T}_pN\)，纤维配对为固定的 \(\eta_p\)。选一个辛线性同构 \((\mathbb R^{2r},\omega_{\mathrm{std}})\to(\mathrm{T}_pN,\eta_p)\)，定理即把它的邻域与 \(S\times\{0\}\subset(S\times\mathbb R^{2r},\sigma+\omega_{\mathrm{std}})\) 的邻域辛同胚地识别。若使用 \(p\) 附近的 Darboux 图，便能直接取这张图与 \(\operatorname{id}_S\) 的乘积；这也展示了法丛平凡时最简单的模型。
\end{example}

\subsection{Lagrange 邻域定理}
\label{g4:subsec:0803:04}

对 Lagrange 子流形 \(L\subset(M,\omega)\)，有 \((\mathrm{T}L)^\omega=\mathrm{T}L\)，辛正交补不再提供横向方向。取代它的是切向与法向之间的非退化配对。按本书的余切辛形式符号，法商与余切丛的识别写成
\begin{equation}\label{g4:eq:0803:lagrangian-normal-pairing}
	\mathrm{T}M|_L/\mathrm{T}L\longrightarrow\mathrm{T}^*L,
	\qquad [v]\longmapsto\bigl(u\mapsto\omega(u,v)\bigr).
\end{equation}
若更换 \(v\) 的切向代表，配对不变；若右边为零，则 \(v\in(\mathrm{T}L)^\omega=\mathrm{T}L\)。维数相同，所以它是光滑丛同构。这里配对把切向量放在第一个变量，符号将与 \(\omega_{\mathrm{can}}=-\mathrm{d}\theta\) 一致。

\begin{lemma}\label{g4:lem:0803:lagrangian-linear-model}
	设 \(L^n\subset(M^{2n},\omega)\) 为无边界 Lagrange 子流形，\(J\) 为 \(\omega\)-相容几乎复结构，\(g=\omega(\cdot,J\cdot)\)，\(h=g|_{\mathrm{T}L}\)。将度量同构记为 \(h^\sharp:\mathrm{T}^*L\to\mathrm{T}L\)。则
	\begin{equation}\label{g4:eq:0803:lagrangian-linear-isomorphism}
		F:\mathrm{T}L\oplus\mathrm{T}^*L\longrightarrow\mathrm{T}M|_L,
		\qquad F(u,p)=u+Jh^\sharp p
	\end{equation}
	是光滑辛丛同构，其中源上的形式为
	\[
		\omega_{\mathrm{can}}\bigl((u,p),(v,r)\bigr)=r(u)-p(v).
	\]
	它在 \(\mathrm{T}L\) 上为恒等，并实现式\eqref{g4:eq:0803:lagrangian-normal-pairing}的逆向法丛识别。
\end{lemma}
\begin{proof}
	先说明 \(J\mathrm{T}L\) 为补丛。若 \(Jw\in\mathrm{T}L\) 且 \(w\in\mathrm{T}L\)，则 Lagrange 条件给出 \(\omega(w,Jw)=0\)；相容正性迫使 \(w=0\)。所以 \(\mathrm{T}L\cap J\mathrm{T}L=0\)，维数相加得到 \(\mathrm{T}M|_L=\mathrm{T}L\oplus J\mathrm{T}L\)。此外 \(\omega(Ju,Jv)=\omega(u,v)=0\) 对切于 \(L\) 的 \(u,v\) 成立，故补丛也是 Lagrange 的。

	因此式\eqref{g4:eq:0803:lagrangian-linear-isomorphism}为光滑线性同构。计算四项配对，两个同类项为零，而混合项为
	\[
		\omega(u,Jh^\sharp r)=h(u,h^\sharp r)=r(u),\qquad
		\omega(Jh^\sharp p,v)=-p(v).
	\]
	故 \(F^*\omega=\omega_{\mathrm{can}}\)。同一计算说明 \([Jh^\sharp p]\) 在式\eqref{g4:eq:0803:lagrangian-normal-pairing}下恰映为 \(p\)。
\end{proof}

\begin{theorem}[Lagrange 邻域定理]\label{g4:thm:0803:lagrangian-neighborhood}
	设 \((M^{2n},\omega)\) 为无边界辛流形，\(i:L^n\hookrightarrow M\) 为闭的无边界 Lagrange 嵌入，\(s:L\to\mathrm{T}^*L\) 为零截面。则存在零截面的开邻域 \(U\subset\mathrm{T}^*L\)、\(L\) 的开邻域 \(V\subset M\)，以及微分同胚 \(\Psi:U\to V\)，使
	\[
		\Psi\circ s=i,\qquad \Psi^*\omega=\omega_{\mathrm{can}}=-\mathrm{d}\theta.
	\]
	若先选定一个相容 \(J\)，还可令 \(\mathrm{d}\Psi\) 沿零截面等于式\eqref{g4:eq:0803:lagrangian-linear-isomorphism}的 \(F\)。不要求 \(L\) 紧致。
\end{theorem}
\begin{proof}
	由\cref{g4:thm:0703:bundle-existence}选取相容几乎复结构 \(J\)。沿零截面有规范分裂
	\[
		s^*\mathrm{T}(\mathrm{T}^*L)=\mathrm{T}L\oplus\mathrm{T}^*L,
	\]
	分别来自零截面的切向量与竖直纤维向量。在余切坐标中，\(\omega_{\mathrm{can}}=\sum_j\mathrm{d}q^j\wedge\mathrm{d}p_j\)，故该分裂中的配对正是 \(r(u)-p(v)\)。\cref{g4:lem:0803:lagrangian-linear-model}于是给出覆盖 \(i\) 的整个环境切丛辛同构 \(F\)，在切向上等于 \(\mathrm{d}i\)。

	由\cref{g4:lem:0803:realize-first-jet}，先取得邻域微分同胚 \(\Phi\)，使 \(\Phi\circ s=i\)、\(\mathrm{d}\Phi|_{s(L)}=F\)。于是 \(\Phi^*\omega\) 与 \(\omega_{\mathrm{can}}\) 沿零截面在整个环境切空间上相等。\cref{g4:thm:0803:relative-moser}给出一个固定零截面并保持一阶恒等的 \(\varphi\)，满足
	\[
		\varphi^*(\Phi^*\omega)=\omega_{\mathrm{can}}.
	\]
	令 \(\Psi=\Phi\circ\varphi\)，便有 \(\Psi^*\omega=\omega_{\mathrm{can}}\)、\(\Psi\circ s=i\) 及 \(\mathrm{d}\Psi|_{s(L)}=F\)。这条证明中的管状邻域与流均允许沿非紧零截面改变宽度，故同样适用于非紧 \(L\)。
\end{proof}

原流形中的 Lagrange 子流形至此被化为余切丛的零截面。辅助 \(J\) 用来选取一张具体的一阶识别，邻域模型的存在则只取决于 Lagrange 条件。文献中通常将这一结论称为\term[Weinsteinguanzhuanglinyudingli]{Weinstein 管状邻域定理}[Weinstein tubular neighborhood theorem]。\footnote{Alan David Weinstein（1943年生），美国数学家，研究辛几何、Poisson 几何与几何力学。}\cite[第2.2--2.3节，定理2.8、2.10，第95--97页]{Cannas2006Symplectic}

\begin{example}\label{g4:exm:0803:circle-neighborhood}
	在 \((\mathbb R^2,\mathrm{d}x\wedge\mathrm{d}y)\) 中取半径 \(R>0\) 的圆周 \(L\)。一维子流形自动为 Lagrange。令 \(\theta\in\mathbb R/2\pi\mathbb Z\) 为圆周坐标，\(p\) 为 \(\mathrm{T}^*L\) 中相对于 \(\mathrm{d}\theta\) 的纤维坐标。对 \(|p|<R^2/4\)，定义
	\[
		\Psi(\theta,p)=\bigl(\sqrt{R^2-2p}\cos\theta,\sqrt{R^2-2p}\sin\theta\bigr).
	\]
	其像为一个圆环，\(p=0\) 映到原圆周。写 \(r=\sqrt{R^2-2p}\)，则 \(r\,\mathrm{d}r=-\mathrm{d}p\)，所以
	\[
		\Psi^*(\mathrm{d}x\wedge\mathrm{d}y)
		=r\,\mathrm{d}r\wedge\mathrm{d}\theta
		=\mathrm{d}\theta\wedge\mathrm{d}p
		=\omega_{\mathrm{can}}.
	\]
	这是一张显式的 Lagrange 邻域同胚。正的 \(p\) 指向圆周内侧，来自本书的辛形式符号；\(p=(R^2-r^2)/2\) 则说明余切变量记录的是径向扫过的有向面积尺度，而不是径向距离本身。
\end{example}

本节的两类邻域定理有相同的证明机制：先实现整个环境切空间上的辛线性同构，再用相对原始与流修正邻域中的形式，并保留已匹配的一阶数据。辛子流形的横向数据仍是一份辛向量丛；Lagrange 子流形的横向数据由切向配对确定为余切丛。这一区别决定了两个模型的形式。
% 纲要条目排为编号小节标题，并进入目录与书签；说明段仍保留于纲要注释中。

% !TeX root = ../../Book4.tex
% 纲要标题：### 8.4 图像与生成函数
% 纲要位置：几何学计划纲要.md，第 3800 行。
% 正文状态：初稿；尚未正式审核。

% 纲要内容开始
% 
% * 闭一形式和局部恰当性
% * 辛映射图像的 Lagrange 性
% * 生成函数及不同投影的适用范围
% 
% 纲要内容结束

\section{图像与生成函数}
\label{g4:sec:0804}

邻域定理把 Lagrange 子流形附近的几何搬到余切丛。在这个模型中，只要某个子流形能投影成一张图，它就由一形式描述；Lagrange 条件使该一形式闭，局部又成为函数的微分。辛映射本身也能通过它的图像进入这个描述。于是，寻找辛映射可以转化为寻找函数，但必须先确认所选投影真的能作为坐标。

\subsection{闭一形式和局部恰当性}
\label{g4:subsec:0804:01}

设 \(L\) 是 \((\mathrm{T}^*Q,-\mathrm{d}\theta)\) 中的 Lagrange 子流形，\(\pi:\mathrm{T}^*Q\to Q\) 为投影。对一点 \(z\in L\)，若 \(\mathrm{d}(\pi|_L)_z\) 可逆，第一卷定理11.3.3的逆函数定理给出 \(z\) 的一个邻域，使其中的 \(L\) 可以写成某个开集 \(U\subset Q\) 上的一形式图像。这里可逆性等价于
\[
	\mathrm{T}_zL\cap\ker\mathrm{d}\pi_z=0.
\]
也就是说，切空间不能含有竖直方向。

\begin{proposition}\label{g4:prop:0804:local-graph}
	设 \(L\subset\mathrm{T}^*Q\) 为嵌入 Lagrange 子流形，且余切投影在 \(z\in L\) 处的限制微分可逆。则缩小邻域后存在光滑实函数 \(S:U\to\mathbb R\)，使
	\[
		L=\{(q,\mathrm{d}S_q):q\in U\}
	\]
	在该邻域内成立。若 \(U\) 连通，则 \(S\) 在加常数的意义下唯一。
\end{proposition}
\begin{proof}
	先由逆函数定理把 \(L\) 写为一形式 \(\alpha\) 的图。\cref{g4:prop:0702:lagrangian-graph}给出图像上的拉回 \(-\mathrm{d}\alpha\)，所以 \(\mathrm{d}\alpha=0\)。把 \(U\) 再缩为星形坐标球，第七章给出的径向积分
	\[
		S(q)=\int_0^1\alpha_{tq}(q)\,\mathrm{d}t
	\]
	满足 \(\mathrm{d}S=\alpha\)。若两个函数给出同一图，它们的微分相同，差在连通域上为常数；反过来加常数不改变图像。
\end{proof}

若 \(L\) 已是整个 \(Q\) 上的闭一形式 \(\alpha\) 的图，局部函数不一定拼成全局函数。沿任意分段光滑闭路 \(\gamma\) 的周期 \(\int_\gamma\alpha\) 必须为零，才可能有 \(\alpha=\mathrm{d}S\)。这个条件也充分：在每个连通分支固定起点，沿道路积分定义 \(S(q)\)；两条道路之差给出闭路，故结果与道路无关。局部用刚才的势函数比较，说明 \(S\) 光滑且微分为 \(\alpha\)。第七章圆周上的 \(c\,\mathrm{d}t\) 图像正好展示了非零周期的障碍。

\begin{example}\label{g4:exm:0804:folded-projection}
	在 \(\mathbb R^2_{q,p}\) 中，曲线 \((q,p)=(s^2,s)\) 是光滑嵌入，并且因维数为一而是 Lagrange 子流形。其投影导数为 \(2s\)，在 \(s=0\) 处为零，所以原点附近的整条曲线不能写成一个 \(q\) 函数的微分图像。在 \(q>0\) 时，两支分别满足
	\[
		p=\sqrt q=\frac{\mathrm{d}}{\mathrm{d}q}\left(\frac23q^{3/2}\right),
		\qquad
		p=-\sqrt q=\frac{\mathrm{d}}{\mathrm{d}q}\left(-\frac23q^{3/2}\right).
	\]
	曲线本身在原点完全光滑；失效的是选 \(q\) 作参数的方式。换用 \(p\) 时，\(q=p^2\) 又是一张光滑图。
\end{example}

% !TeX root = ../Book4.tex
\begin{figure}[htbp]
	\centering
	\begin{tikzpicture}[x=1.5cm,y=1.2cm,>=Stealth,font=\small]
		\draw[->] (-0.45,0)--(2.7,0) node[right] {\(q\)};
		\draw[->] (0,-1.6)--(0,1.65) node[above] {\(p\)};
		\draw[GeometryBlue,thick,domain=-1.35:1.35,samples=80,smooth,variable=\t] plot ({\t*\t},\t);
		\draw[dashed,GeometryGray] (1,-1)--(1,1);
		\draw[->,GeometryBlue] (1,0.84)--(1,0.12);
		\draw[->,GeometryBlue] (1,-0.84)--(1,-0.12);
		\fill[GeometryBlue] (1,1) circle (2pt) node[above left] {\((1,1)\)};
		\fill[GeometryBlue] (1,-1) circle (2pt) node[below left] {\((1,-1)\)};
		\fill (1,0) circle (1.6pt) node[below right] {\(1\)};
		\fill (0,0) circle (1.6pt) node[below left] {\(0\)};
		\node[right] at (1.9,1.25) {\(q=p^2\)};
	\end{tikzpicture}
	\caption[Lagrange 曲线与底投影]{光滑曲线 \(q=p^2\) 向 \(q\) 轴投影时，两点可落到同一个底点，原点处的切线还完全竖直。曲线的光滑性并不保证它是关于 \(q\) 的单值图像。}
	\label{g4:fig:0804:graph-projection}
\end{figure}


\subsection{辛映射图像的 Lagrange 性}
\label{g4:subsec:0804:02}

设 \((M,\omega_M)\)、\((N,\omega_N)\) 都有实维数 \(2n\)。在乘积上取
\begin{equation}\label{g4:eq:0804:twisted-product}
	\Omega=-\pi_M^*\omega_M+\pi_N^*\omega_N.
\end{equation}
它在两个因子上分别非退化，因而是辛形式。记 \(M^-=(M,-\omega_M)\)，这个辛乘积也写成 \(M^-\times N\)。负号的作用是使同一向量在两因子中的配对互相抵消。

\begin{proposition}\label{g4:prop:0804:symplectic-graph}
	设 \(F:M\to N\) 为同维辛流形间的光滑映射。其图像 \(\Gamma_F\subset M^-\times N\) 为 Lagrange 子流形，当且仅当 \(F^*\omega_N=\omega_M\)。在这个条件下，\(F\) 是局部辛同胚；若 \(F\) 为双射，则是全局辛同胚。
\end{proposition}
\begin{proof}
	映射 \(j_F(x)=(x,F(x))\) 是嵌入，其像的维数为 \(2n\)，即乘积维数的一半。由两张投影的复合关系，
	\[
		j_F^*\Omega=-\omega_M+F^*\omega_N.
	\]
	因此图像为 Lagrange 子流形恰好等价于所述拉回等式。若 \(\mathrm{d}F_xv=0\)，则对一切 \(w\) 有 \(\omega_M(v,w)=0\)，故 \(v=0\)。同维性使 \(\mathrm{d}F_x\) 可逆，所以 \(F\) 局部为微分同胚。双射时，各局部光滑逆拼成全局光滑逆。
\end{proof}

特别地，对角线 \(\Delta\subset M^-\times M\) 为 Lagrange 子流形。\cref{g4:thm:0803:lagrangian-neighborhood}把其邻域与 \(\mathrm{T}^*M\) 中零截面的邻域辛同构。选定这样的邻域对应后，靠近恒等映射的辛同胚便可以用靠近零的一形式来描述。下面说明这里为什么采用 \(C^1\) 接近，而不只要求点的位置接近。\cite[第2.4节，第98--99页]{Cannas2006Symplectic}

\begin{proposition}\label{g4:prop:0804:near-identity-forms}
	设 \((M,\omega)\) 为闭连通辛流形，固定一个把零截面送到对角线的邻域辛同胚
	\[
		\Psi:U\subset\mathrm{T}^*M\longrightarrow V\subset M^-\times M,
		\qquad \Psi(x,0)=(x,x).
	\]
	充分 \(C^1\) 接近恒等映射的辛同胚，与充分 \(C^1\) 接近零的光滑闭一形式局部一一对应；对应由 \(\Psi(\operatorname{graph}\alpha)=\Gamma_F\) 给出。该对应下，\(F\) 的固定点恰好对应 \(\alpha\) 的零点。
\end{proposition}
\begin{proof}
	先说明紧性怎样保证投影可逆。在闭流形上，充分 \(C^1\) 接近恒等映射的光滑映射 \(a:M\to M\) 必为微分同胚。取有限组坐标球，其中较小同心球覆盖 \(M\)。对落入同一较小球且足够接近的两个点，连接坐标线段位于较大球中；若坐标表达满足 \(\|\mathrm{D}a-I\|<1/2\)，沿线段积分给出
	\[
		|a(x)-a(y)|\ge\tfrac12|x-y|.
	\]
	由有限覆盖，可选统一的接近尺度。距离不小于该尺度的点对，则由 \(a\) 与恒等映射充分 \(C^0\) 接近而保持不同。于是 \(a\) 单射，其微分也处处可逆。像既开，又因紧性而闭；连通性给出满射，局部光滑逆遂成为全局光滑逆。

	若 \(F\) 接近恒等映射，紧性保证其图包含在 \(V\) 中。令 \(\sigma=\Psi^{-1}\circ j_F\)、\(a=\pi\circ\sigma\)，则 \(\sigma\) 接近零截面，\(a\) 接近恒等映射，所以 \(a\) 为微分同胚。重参数化 \(\sigma\circ a^{-1}\) 是一个真正的余切截面，记为 \(\alpha\)。图像经辛同胚变成 Lagrange 子流形，因此 \(\alpha\) 闭。

	反过来，对充分小的闭一形式 \(\alpha\)，写 \(\Psi\circ s_\alpha=(a_\alpha,b_\alpha)\)。两张映射都 \(C^1\) 接近恒等映射，故均为微分同胚；取 \(F=b_\alpha\circ a_\alpha^{-1}\)，所得子流形就是 \(F\) 的图。它是 Lagrange 子流形，所以 \(F\) 辛。两次构造都以同一个子流形为中介，彼此互逆。图像与对角线相交，当且仅当 \(\alpha\) 的图与零截面相交；交点 \((x,0)\) 被 \(\Psi\) 送到 \((x,x)\)，恰好保留所述固定点对应。
\end{proof}

\begin{corollary}\label{g4:cor:0804:near-identity-fixed-points}
	设 \((M^{2n},\omega)\) 为闭连通辛流形，\(n\ge1\)，且 \(H^1_{\mathrm{dR}}(M)=0\)。则充分 \(C^1\) 接近恒等映射的辛同胚至少有两个固定点。
\end{corollary}
\begin{proof}
	前一命题给出闭一形式 \(\alpha\)。由上同调假设，\(\alpha=\mathrm{d}S\)。若 \(S\) 非常值，它在紧流形上的最大值点与最小值点不同，且两点的微分都为零；若 \(S\) 常值，则 \(\alpha=0\)，对应映射为恒等映射。两种情形均给出至少两个固定点。
\end{proof}

这条结论把一个局部辛几何问题变成了函数的极值问题。对一般辛同胚，其图可能离开所选对角线邻域；即使仍在邻域内，也要另查向零截面的投影是否可逆。

\subsection{生成函数及不同投影的适用范围}
\label{g4:subsec:0804:03}

在源、靶余切坐标分别写 \((q,p)\)、\((Q,P)\)。式\eqref{g4:eq:0804:twisted-product}成为
\[
	\Omega=-\sum_i\mathrm{d}q^i\wedge\mathrm{d}p_i
	+\sum_i\mathrm{d}Q^i\wedge\mathrm{d}P_i
	=\mathrm{d}\lambda,
	\qquad \lambda=\sum_i p_i\,\mathrm{d}q^i-P_i\,\mathrm{d}Q^i.
\]
因此在辛映射的图上，\(\lambda\) 闭。若 \((q,Q)\) 是这张图上的局部坐标，就可以把它的局部原始函数写为 \(S(q,Q)\)。所得方程是
\begin{equation}\label{g4:eq:0804:position-generating}
	p_i=\frac{\partial S}{\partial q^i},
	\qquad P_i=-\frac{\partial S}{\partial Q^i}.
\end{equation}
这就是以两组位置为变量的生成函数。关键还在于：给定 \((q,p)\) 后，能否从第一组方程解出 \(Q\)。\cite[第2.5节，第99--100页]{Cannas2006Symplectic}

\begin{proposition}\label{g4:prop:0804:generating-function}
	设 \(S\) 为 \(\mathbb R_q^n\times\mathbb R_Q^n\) 某开集上的光滑实函数。若在 \((q_0,Q_0)\) 处，混合二阶导数矩阵
	\[
		S_{qQ}=\left(\frac{\partial^2S}{\partial q^i\partial Q^j}\right)_{i,j=1}^n
	\]
	可逆，则式\eqref{g4:eq:0804:position-generating}在对应点附近确定一个局部辛同胚 \((q,p)\mapsto(Q,P)\)。反过来，若辛同胚的图能以 \((q,Q)\) 为局部坐标，则在充分小的邻域内存在这样的 \(S\)，且 \(S_{qQ}\) 必可逆。
\end{proposition}
\begin{proof}
	对方程 \(p=S_q(q,Q)\)，关于 \(Q\) 的导数是可逆矩阵 \(S_{qQ}\)。第一卷定理11.3.5的带参数隐函数定理给出唯一光滑局部解 \(Q=Q(q,p)\)，再以 \(P=-S_Q(q,Q(q,p))\) 定义输出。相应图上有 \(\lambda=\mathrm{d}S\)，故 \(\Omega\) 的拉回为零；\cref{g4:prop:0804:symplectic-graph}证明所得映射为局部辛同胚。

	反向，\(\lambda\) 在图上闭，以 \((q,Q)\) 为坐标并缩小到球后，存在原始函数 \(S\)。比较 \(\mathrm{d}q\)、\(\mathrm{d}Q\) 系数即得两组方程。图向源空间 \((q,p)\) 的投影本来是局部微分同胚，而在所选参数下它为 \((q,Q)\mapsto(q,S_q)\)，导数的对角块为 \(I,S_{qQ}\)。因此混合矩阵必须可逆。
\end{proof}

恒等映射的图满足 \(Q=q\)，所以 \((q,Q)\) 只有 \(n\) 个独立变量，不能用于这一描述。换成 \((q,P)\) 则恰好适用。因为
\[
	\lambda+\mathrm{d}(Q\cdot P)
	=\sum_i p_i\,\mathrm{d}q^i+Q^i\,\mathrm{d}P_i,
\]
当 \((q,P)\) 是图上的局部坐标时，右侧有局部原始函数 \(F(q,P)\)，满足
\begin{equation}\label{g4:eq:0804:mixed-generating}
	p=F_q(q,P),\qquad Q=F_P(q,P).
\end{equation}
给定 \(F\) 时，条件 \(\det F_{qP}\ne0\) 同样保证先从第一式解出 \(P=P(q,p)\)，再定义 \(Q\)，得到局部辛同胚。证明与上一个命题相同：隐函数定理负责解方程，微分 \(p\,\mathrm{d}q+Q\,\mathrm{d}P=\mathrm{d}F\) 则负责检验辛条件。在两种参数同时可用的区域里，\(F=S+Q\cdot P\)，其中要先用 \(P=-S_Q\) 把 \(Q\) 写成 \((q,P)\) 的函数。

\begin{example}\label{g4:exm:0804:two-projections}
	取 \(S(q,Q)=-q\cdot Q\)，便有 \(p=-Q\)、\(P=q\)，得到标准四分之一周旋转 \((q,p)\mapsto(-p,q)\)。它的 \((q,Q)\) 投影可逆，而 \((q,P)=(q,q)\) 投影失去一半秩。

	取 \(F(q,P)=q\cdot P-V(q)\)，其中 \(V\) 为任意光滑函数，则 \(p=P-\mathrm{d}V_q\)、\(Q=q\)，得到 \((q,p)\mapsto(q,p+\mathrm{d}V_q)\)。\(F_{qP}=I\)，故在整个坐标域上都可解。\(V=0\) 时就是恒等映射。两种投影分别适合不同的图像，选错投影会使生成函数的表达失效，而辛映射本身仍可光滑。
\end{example}

一个局部函数由混合非退化条件产生的是局部辛同胚。把这些公式用于整个开集时，还须检查反解能否单值延拓以及映射是否单射；逐点非零的混合行列式不负责这些整体问题。下一节将把局部坐标的自由与整体嵌入的限制并列比较。
% 纲要条目排为编号小节标题，并进入目录与书签；说明段仍保留于纲要注释中。

% !TeX root = ../../Book4.tex
% 纲要标题：### 8.5 局部自由与整体障碍
% 纲要位置：几何学计划纲要.md，第 3806 行。
% 正文状态：初稿；尚未正式审核。

% 纲要内容开始
% 
% 在这里准确提出球—柱嵌入问题并预告非挤压；证明安排在第24章，不作为本章标准形的前提。
% 
% ---
% 
% 纲要内容结束

\section{局部自由与整体障碍}
\label{g4:sec:0805}

Darboux 定理保证每一点附近都有标准辛坐标。它允许我们把足够小的球放进坐标邻域，却没有规定这些坐标能覆盖多大的区域。若给定一个有固定半径的球，要求把整个球放进另一个辛流形，问题便会涉及坐标之外的整体约束。

\ProseParagraphBreak

固定标准辛空间 \((\mathbb R^{2n},\omega_0)\)，坐标按 \((q^1,p_1,\ldots,q^n,p_n)\) 成对排列，\(\omega_0=\sum_j\mathrm{d}q^j\wedge\mathrm{d}p_j\)。对 \(R,r>0\)，定义开球与开辛柱
\begin{equation}\label{g4:eq:0805:ball-cylinder}
	\begin{aligned}
		B^{2n}(R)&=\left\{\sum_{j=1}^n\bigl((q^j)^2+p_j^2\bigr)<R^2\right\},\\
		Z^{2n}(r)&=\{(q^1)^2+p_1^2<r^2\}
		=B^2(r)\times\mathbb R^{2n-2}.
	\end{aligned}
\end{equation}
这里柱只约束第一对辛共轭方向，其余方向可以无限延伸。\term[xinqianru]{辛嵌入}[symplectic embedding]是保持辛形式的光滑嵌入；等维时它到其像为辛同胚。球—柱问题问的是：何时存在
\[
	\varphi:B^{2n}(R)\hookrightarrow Z^{2n}(r),
	\qquad \varphi^*\omega_0=\omega_0?
\]

\begin{theorem}[非挤压定理]\label{g4:thm:0805:nonsqueezing-preview}
	设 \(n\ge1\)、\(R,r>0\)。标准辛球 \(B^{2n}(R)\) 辛嵌入标准辛柱 \(Z^{2n}(r)\) 的充要条件为 \(R\le r\)。
\end{theorem}

其中 \(R\le r\) 时，恒等包含已经给出嵌入。反向的结论称为 \term[feijiya]{Gromov 非挤压定理}[Gromov nonsqueezing theorem]\footnote{Mikhael Leonidovich Gromov（1943年生），俄裔法国数学家，研究几何、拓扑及几何分析。}，可参见 Cannas da Silva 在伪全纯曲线应用中的陈述。\cite[第3.6节，第124页]{Cannas2006Symplectic} 二维情形只需面积：辛嵌入保持面积，故 \(\pi R^2\le\pi r^2\)。对于 \(n\ge2\)，\chapref{g4:ch:24}将在 \(S^2\times\mathrm{T}^{2n-2}\) 中使用指定球面类、求值映射的度与局部面积估计，完成必要性的证明。下面先证明线性情形，并具体解释高维结论为什么超出体积比较。

\begin{proposition}\label{g4:prop:0805:linear-nonsqueezing}
	设 \(A\) 为 \(\mathbb R^{2n}\) 上的线性辛同构，且 \(A(B^{2n}(R))\subset Z^{2n}(r)\)。则 \(R\le r\)。
\end{proposition}
\begin{proof}
	在本节成对排列的坐标中，令 \(\Omega=\operatorname{diag}(J_2,\ldots,J_2)\)，其中 \(J_2=\begin{psmallmatrix}0&1\\-1&0\end{psmallmatrix}\)。把 \(Ax\) 的第一对分量写成 \(a\cdot x\)、\(b\cdot x\)，其中 \(a,b\in\mathbb R^{2n}\) 由矩阵相应两行的系数组成。由 \(A^{\mathsf T}\Omega A=\Omega\) 可得 \(A\Omega A^{\mathsf T}=\Omega\)，所以
	\[
		a^{\mathsf T}\Omega b=1.
	\]
	标准交替矩阵 \(\Omega\) 正交，故 \(1\le |a||b|\)。另一方面，球像包含于柱，分别取上确界得到
	\[
		R|a|=\sup_{|x|<R}|a\cdot x|\le r,
		\qquad R|b|\le r.
	\]
	两式相乘可得 \(R^2\le R^2|a||b|\le r^2\)，所以 \(R\le r\)。
\end{proof}

一般非线性辛嵌入的微分在每点都满足相同的线性关系，但微分随位置改变，无法直接把这条点态估计积分成固定球像的半径下界。这正是完整非挤压证明需要新方法的地方。

\ProseParagraphBreak

体积能给出另一种约束。辛映射保持 \(\omega_0^n/n!\)，而球的辛体积为 \(\pi^nR^{2n}/n!\)。因此把球嵌入一个给定有限体积区域时，目标的体积必须至少这么大。可是 \(n\ge2\) 时每个 \(Z^{2n}(r)\) 都有无限体积，这个必要条件对任何 \(r>0\) 都成立。即使要求严格保持体积，也确实能把球挤进很细的柱：对 \(0<\varepsilon<1\)，取
\[
	A_\varepsilon(q^1,p_1,q^2,p_2,\ldots)
	=(\varepsilon q^1,\varepsilon p_1,
	\varepsilon^{-1}q^2,\varepsilon^{-1}p_2,\ldots).
\]
它的行列式为 \(1\)，并将 \(B^{2n}(R)\) 送入 \(Z^{2n}(\varepsilon R)\)。但其拉回为
\[
	A_\varepsilon^*\omega_0
	=\varepsilon^2\mathrm{d}q^1\wedge\mathrm{d}p_1
	+\varepsilon^{-2}\mathrm{d}q^2\wedge\mathrm{d}p_2
	+\sum_{j=3}^n\mathrm{d}q^j\wedge\mathrm{d}p_j,
\]
所以它没有保持每一对方向的辛配对。体积记录全部方向共同产生的乘积，非挤压却还能检测指定辛二维方向上的面积尺度。

\ProseParagraphBreak

柱中受约束的二维平面也必须说明。如果改成同时限制两个位置方向，即
\[
	C(r)=\{(q^1)^2+(q^2)^2<r^2\},\qquad n\ge2,
\]
则辛线性变换
\[
	(q^1,q^2,p_1,p_2)\longmapsto
	(\varepsilon q^1,\varepsilon q^2,
	\varepsilon^{-1}p_1,\varepsilon^{-1}p_2)
\]
在其他方向上保持恒等，便能把 \(B^{2n}(R)\) 送入 \(C(\varepsilon R)\)。这里两个位置方向构成各向同性平面，缩小它们时可以分别放大各自配对的动量方向。式\eqref{g4:eq:0805:ball-cylinder}同时约束的却是互相配对的 \(q^1,p_1\)。辛几何的局部标准形与这种整体区别相容：局部总能选坐标，给定大小的区域能否整体嵌入，还要保留配对所携带的信息。
% 本节沿用纲要的连续叙述安排，不另造小节标题。


\begin{chaptersummary}
Moser 方程 \(\iota_{X_t}\omega_t=-\alpha_t\) 把形式的变化转成坐标运动。只要 \(\dot\omega_t=\mathrm{d}\alpha_t\)、\(\omega_t\) 始终非退化，而且流在所需时间内存在，拉回求导就给出 \(\phi_t^*\omega_t=\omega_0\)。闭流形上的上同调类不变提供光滑原始族，紧致性保证整个时间区间上的同位。

\ProseParagraphBreak

Darboux 定理先用辛基匹配一点处的形式，再在共同小球内沿直线插值。径向原始在原点二阶消失，使 Moser 向量场也二阶消失；轨道估计保证局部流存在到时间1，变分方程保证指定的一阶标架不变。由此，同维辛形式在各点附近都具有相同的标准形。

\ProseParagraphBreak

沿子流形作同样的修正，需要区分切向拉回与整个环境张量。相对同伦提供沿子流形消失的原始；匹配环境线性数据后，原始的一阶导数也消失，流便同时固定子流形和环境切空间上的微分。辛子流形的邻域由自身辛形式及辛法丛决定，Lagrange 子流形的邻域则由其余切丛给出。

\ProseParagraphBreak

辛同胚的图像在反号乘积中是 Lagrange 子流形。选择合适的投影后，闭一形式成为函数的微分，产生生成函数；投影退化时，需要换用另一组变量。局部标准形并未消除整体限制：非挤压约束一个共轭坐标平面上的投影尺度，而高维体积本身无法给出这个约束。
\end{chaptersummary}

\BookExercises
\begin{exercise}\label{g4:ex:0801:sphere-area}
	设 \(\omega_0,\omega_1\) 是定向球面 \(S^2\) 上的两个正面积形式。证明：存在从恒等映射出发的光滑同位 \(\phi_t:S^2\to S^2\)，使 \(\phi_1^*\omega_1=\omega_0\)，当且仅当 \(\int_{S^2}\omega_0=\int_{S^2}\omega_1\)。说明正性在直线路径中起什么作用。
\end{exercise}
\begin{solution}
	同位中的微分同胚保持定向，故换元公式给出积分相等。反之，第二卷第7.4.2小节的球面上同调计算说明，\(S^2\) 上闭二形式的上同调类由积分决定。因此积分相等给出 \([\omega_0]=[\omega_1]\)。取 \(\omega_t=(1-t)\omega_0+t\omega_1\)。在每一正定向基上，两个端点形式的值都为正，凸组合仍为正，所以所有 \(\omega_t\) 都非退化；二维中二形式自动闭。整条路径的上同调类不变，\cref{g4:thm:0801:moser}给出所需同位。这里正性保证两形式相对于同一定向具有同号系数，从而直线插值不会经过零。
\end{solution}

\begin{exercise}\label{g4:ex:0801:finite-equivariance}
	设有限群 \(G\) 光滑作用于闭流形 \(M\)，并保持光滑辛形式族 \(\omega_t\)、\(0\le t\le1\) 中的每个形式。若 \([\omega_t]\) 不变，证明可选 Moser 同位 \(\phi_t\) 与每个群元素的作用交换。
\end{exercise}
\begin{solution}
	由光滑原始族的构造，取 \(\mathrm{d}\alpha_t=\dot\omega_t\)，再令
	\[
		\overline\alpha_t=\frac1{|G|}\sum_{g\in G}g^*\alpha_t.
	\]
	有限求和保持光滑性；因 \(g^*\omega_t=\omega_t\)，有 \(\mathrm{d}\overline\alpha_t=\dot\omega_t\)，且 \(\overline\alpha_t\) 在群作用下不变。用 \(\iota_{X_t}\omega_t=-\overline\alpha_t\) 定义向量场。对 \(v\in\mathrm{T}_{g x}M\)，形式的不变性给出
	\[
		\omega_{t,gx}(\mathrm{d}g_xX_t(x),v)
		=\omega_{t,x}(X_t(x),(\mathrm{d}g_x)^{-1}v)
		=-\overline\alpha_{t,gx}(v).
	\]
	解的唯一性说明 \(X_t(gx)=\mathrm{d}g_xX_t(x)\)。因此 \(g\phi_t(x)\) 与 \(\phi_t(gx)\) 满足同一初值问题，由常微分方程唯一性相等。\(M\) 的紧致性保证流存在于全部时间区间，通常的拉回计算给出 \(\phi_t^*\omega_t=\omega_0\)。
\end{solution}

\begin{exercise}\label{g4:ex:0801:nonlinear-path}
	在 \(\mathbb R^4_{q^1,q^2,p_1,p_2}\) 上，\(\omega_0\) 与 \(-\omega_0\) 之间的直线路径经过零形式。证明它们仍可由一条处处非退化的光滑二形式路径连接，并显式写出相应的 Moser 同位。提示：只旋转 \((p_1,p_2)\) 平面。
\end{exercise}
\begin{solution}
	令 \(R_t=\begin{pmatrix}\cos\pi t&-\sin\pi t\\\sin\pi t&\cos\pi t\end{pmatrix}\)、\(A_t(q,p)=(q,R_tp)\)，并取
	\[
		\omega_t=A_t^*\omega_0
		=\sum_{i,j}(R_t)_{ij}\,\mathrm{d}q^i\wedge\mathrm{d}p_j.
	\]
	每个 \(A_t\) 可逆，故 \(\omega_t\) 始终辛；\(R_0=I\)、\(R_1=-I\)，给出要求的端点。取 \(\phi_t=A_t^{-1}\)，直接得到 \(\phi_t^*\omega_t=\omega_0\)。也可核对 Moser 方程：令
	\[
		\alpha_t=-\sum_i(\dot R_tp)_i\,\mathrm{d}q^i,
		\qquad X_t(q,p)=(0,-R_t^{-1}\dot R_tp).
	\]
	则 \(\mathrm{d}\alpha_t=\dot\omega_t\)，而 \(\iota_{X_t}\omega_t=\sum_i(\dot R_tp)_i\mathrm{d}q^i=-\alpha_t\)。矩阵族 \(R_t\) 与其导数交换，故 \(A_t^{-1}\) 正是这个场的流，并且对所有 \(t\) 全局存在。端点之间存在好路径，并不意味着任意插值路径都可用。
\end{solution}

\begin{exercise}\label{g4:ex:0802:higher-jets}
	设 \(\omega\) 是 \(\mathbb R^{2n}\) 中原点附近的光滑辛形式，且 \(\omega-\omega_0\) 的坐标系数在原点的所有低于 \(m\) 阶的导数都为零，其中 \(m\ge1\)。证明可在一个邻域上选取 \(\phi\)，使
	\[
		\phi^*\omega=\omega_0,\qquad
		\phi(x)=x+O(|x|^{m+1}).
	\]
	因而 \(\phi\) 与恒等映射在原点的所有至 \(m\) 阶导数相同。
\end{exercise}
\begin{solution}
	记 \(\beta=\omega-\omega_0\)。光滑性及题设给出 \(\|\beta_x\|\le C|x|^m\)。径向原始满足
	\[
		|\alpha_x(v)|
		\le\int_0^1sC|sx|^m|x||v|\,\mathrm{d}s
		=\frac C{m+2}|x|^{m+1}|v|.
	\]
	在整条路径共同非退化的小球内，收缩映射的逆有统一界，故 Moser 场满足 \(|X_t(x)|\le C'|x|^{m+1}\)。正文的轨道估计先保证一个固定小球中全部初值的解存在到时间1。为记录对初值大小的依赖，再取足够小的 \(x\ne0\)：只要轨道仍在 \(B_{2|x|}\) 中，就有
	\[
		|\phi_t(x)-x|\le C'2^{m+1}|x|^{m+1}.
	\]
	当 \(C'2^{m+1}|x|^m<1/2\) 时，它不能首次到达该球边界，于是上述估计对 \(0\le t\le1\) 全部成立。取 \(\phi=\phi_1\)，Moser 恒等式给出拉回结论，估计给出所需消失阶数。光滑函数 \(\phi-\operatorname{id}\) 若有阶数不超过 \(m\) 的首个非零齐次 Taylor 项，沿该多项式非零的某条直线便不可能是 \(O(|x|^{m+1})\)。所以这些导数全为零。
\end{solution}

\begin{exercise}\label{g4:ex:0802:radial-disc}
	设 \(D\subset\mathbb R^2\) 为闭单位圆盘，\(f(r)=h(r^2)>0\)，其中 \(h\) 在 \([0,1]\) 的邻域光滑，并且
	\[
		2\int_0^1f(s)s\,\mathrm{d}s=1.
	\]
	构造固定边界每一点的光滑微分同胚 \(F:D\to D\)，使 \(F^*(\mathrm{d}x\wedge\mathrm{d}y)=f(r)\,\mathrm{d}x\wedge\mathrm{d}y\)。若 \(f=1\) 在边界附近成立，进一步证明可使 \(F\) 在边界附近恒等。
\end{exercise}
\begin{solution}
	取
	\[
		R(r)=\left(2\int_0^r f(s)s\,\mathrm{d}s\right)^{1/2},
		\qquad F(x,y)=\frac{R(r)}r(x,y)\quad(r>0),\qquad F(0)=0.
	\]
	虽然极坐标在原点失效，比例因子却有光滑表达式
	\[
		\frac{R(r)}r=\left(\int_0^1h(tr^2)\,\mathrm{d}t\right)^{1/2}.
	\]
	它是 \(r^2=x^2+y^2\) 的正光滑函数，在原点的值为 \(\sqrt{h(0)}\)。故 \(F\) 在原点光滑且微分可逆。对 \(r>0\)，有 \(R(r)R'(r)=f(r)r>0\)，所以 \(R\) 严格递增；归一化条件给出 \(R(1)=1\)。它保持角度，因而是圆盘上的双射，并逐点固定边界。径向和切向导数均非零；在边界附近延拓同一公式后应用逆函数定理，可知逆映射也光滑。

	在穿孔圆盘上，\(\mathrm{d}x\wedge\mathrm{d}y=r\,\mathrm{d}r\wedge\mathrm{d}\theta\)，所以拉回为 \(R(r)R'(r)\,\mathrm{d}r\wedge\mathrm{d}\theta=f(r)r\,\mathrm{d}r\wedge\mathrm{d}\theta\)，由光滑性延伸到原点。若 \(f=1\) 在 \(r\ge r_0\) 成立，则对这些 \(r\)，
	\[
		R(r)^2=1-2\int_r^1s\,\mathrm{d}s=r^2,
	\]
	故 \(F\) 在该边界邻域恒等。归一化条件也正是两个面积形式总面积相等的条件。
\end{solution}

\begin{exercise}\label{g4:ex:0803:normal-obstruction}
	在 \(S^2\) 上固定正面积形式 \(\sigma\)。令 \(M_0=M_1=S^2\times S^2\)，但分别赋予辛形式
	\[
		\omega_0=\pi_1^*\sigma+\pi_2^*\sigma,
		\qquad\omega_1=\tfrac12(\pi_1^*\sigma+\pi_2^*\sigma).
	\]
	固定 \(p\in S^2\)，取 \(S_0=S^2\times\{p\}\)、\(S_1=\{(x,x):x\in S^2\}\)。证明 \(f(x,p)=(x,x)\) 是受限辛形式之间的辛同胚，却不能延伸为环境邻域之间的辛同胚。提示：比较辛法丛的第一陈类。
\end{exercise}
\begin{solution}
	两条嵌入从 \(S^2\) 拉回的形式都等于 \(\sigma\)，故 \(f\) 辛。在 \(S_0\) 上，辛法丛为 \(\{0\}\oplus\mathrm{T}_pS^2\)，是平凡的定向二平面丛。在对角线上，切空间为 \(\{(u,u)\}\)；与所有此类向量辛正交的向量满足两分量之和为零，所以辛法丛为
	\[
		\{(v,-v):v\in\mathrm{T}_xS^2\},\qquad
		\omega_1((v,-v),(w,-w))=\sigma(v,w).
	\]
	因此它作为辛向量丛同构于 \(\mathrm{T}S^2\)。上一章对球面切丛的计算给出其第一陈类为 \(2u\)，其中 \(u\) 是按给定定向归一化的生成元；平凡丛的第一陈类为零。两辛法丛不能同构。任何延伸 \(f\) 的环境辛同胚，其微分都会把切丛送到切丛，并保持辛正交补及其配对，因而诱导两辛法丛同构，矛盾。子流形自身的辛形式并未包含全部邻域数据。
\end{solution}

\begin{exercise}\label{g4:ex:0804:torus-diagonal}
	取 \(M=\mathbb R^2/\mathbb Z^2\)、\(\omega=\mathrm{d}q\wedge\mathrm{d}p\)，用 \((q,p;u,v)\) 表示余切向量 \(u\,\mathrm{d}q+v\,\mathrm{d}p\)。证明在零截面附近，
	\[
		\Psi(q,p;u,v)=\bigl((q+v/2,p-u/2),(q-v/2,p+u/2)\bigr)
	\]
	是到反号乘积中对角线附近的辛同胚，所有底坐标按模 \(\mathbb Z^2\) 理解。求小平移 \(F_{a,b}(q,p)=(q+a,p+b)\) 的图像对应的一形式，并用它检验近恒等不动点推论中的上同调条件。
\end{exercise}
\begin{solution}
	直接展开可得
	\[
		\Psi^*(-\pi_1^*\omega+\pi_2^*\omega)
		=\mathrm{d}q\wedge\mathrm{d}u+\mathrm{d}p\wedge\mathrm{d}v,
	\]
	这正是余切丛的典范辛形式。若限制 \(|u|,|v|<\delta<1/2\)，两点坐标差有唯一落在 \((-\delta,\delta)\) 的提升，由差恢复 \(u,v\)，再由第一点恢复 \(q,p\)。因此公式具有光滑逆，给出所述邻域辛同胚。

	图像条件为第二点减第一点等于 \((a,b)\)，即 \(-v=a\)、\(u=b\)。所以对应的一形式是 \(\alpha=b\,\mathrm{d}q-a\,\mathrm{d}p\)。它闭；沿两个基本坐标圆的周期分别为 \(b\) 和 \(-a\)。对充分小且非零的 \((a,b)\)，它无零点且不恰当，相应平移也没有不动点。第8.4节在 \(H^1_{\mathrm{dR}}(M)=0\) 时把闭一形式写成 \(\mathrm{d}f\)，再用极值给出不动点；本例的非零周期正好阻止了这一步。
\end{solution}

\begin{exercise}\label{g4:ex:0804:transverse-intersections}
	设 \(f:Q\to\mathbb R\) 光滑。证明 \(\operatorname{graph}(\mathrm{d}f)\) 与零截面的交点恰对应于 \(f\) 的临界点；在临界点处，两 Lagrange 子流形横截相交，当且仅当 \(f\) 的二阶导数矩阵可逆。应用于
	\[
		f(q)=\varepsilon\sum_{j=1}^n\cos(2\pi q^j)
		\quad\text{在 }\mathbb T^n\text{ 上},\qquad \varepsilon\ne0,
	\]
	计算交点数，并说明如何把这一构造移入任意给定的紧 Lagrange 环面的邻域。
\end{exercise}
\begin{solution}
	图像落在零截面上当且仅当 \(\mathrm{d}f_q=0\)。在这样的点，余切坐标把两切空间分别写成 \(\{(v,0)\}\) 与 \(\{(v,Hv)\}\)，其中 \(H=(\partial_i\partial_jf)(q)\)。交空间同构于 \(\ker H\)。两者维数各为 \(n\)，环境维数为 \(2n\)，所以切空间之和等于环境空间，当且仅当交空间为零，也就是 \(H\) 可逆。在临界点处，变坐标时含一阶导数的项消失，故这个判别不依赖坐标。

	给定函数的临界点要求每个 \(q^j\) 等于 \(0\) 或 \(1/2\) 模整数，共 \(2^n\) 个。二阶导数矩阵为对角阵，其第 \(j\) 项为 \(-4\pi^2\varepsilon\cos(2\pi q^j)\)，在这些点全不为零，因此所有交点横截。若 \(L\) 是给定的紧 Lagrange 环面，选择 \(L\cong\mathbb T^n\)，并应用邻域定理。\(\mathrm{d}f\) 随 \(\varepsilon\to0\) 一致趋于零，紧致性保证充分小的图像落在选定余切邻域中。把图像送回环境，便得到与 \(L\) 恰有 \(2^n\) 个横截交点的 Lagrange 子流形。
\end{solution}

\begin{exercise}\label{g4:ex:0804:composition}
	设 \(S(q,Q)\)、\(T(Q,R)\) 生成相继的局部辛同胚，采用正文的符号 \(p=S_q\)、\(P=-S_Q=T_Q\)、\(r=-T_R\)。假设在某点 \(S_Q+T_Q=0\)，且
	\[
		H=S_{QQ}+T_{QQ}
	\]
	可逆。证明可以解出 \(Q=Q(q,R)\)，并由
	\[
		U(q,R)=S(q,Q(q,R))+T(Q(q,R),R)
	\]
	生成复合映射。计算 \(U_{qR}\)。对 \(S_t(q,Q)=-|Q-q|^2/(2t)\)、\(t\ne0\)，计算参数 \(t,s\) 的复合，其中 \(s\ne0\)、\(t+s\ne0\)。
\end{exercise}
\begin{solution}
	第一卷定理11.3.5的带参数隐函数定理给出所述光滑函数 \(Q(q,R)\)。链式法则中含 \(\mathrm{d}Q\) 的系数正是 \(S_Q+T_Q=0\)，所以
	\[
		U_q=S_q=p,\qquad U_R=T_R=-r.
	\]
	中间动量的等式也已由驻点方程保证，故这些方程描述复合关系。对驻点方程按 \(R\) 求导，得 \(H Q_R+T_{QR}=0\)，于是
	\[
		U_{qR}=-S_{qQ}H^{-1}T_{QR}.
	\]
	两原生成函数的混合矩阵都可逆，\(H\) 也可逆，因此复合的混合矩阵可逆，确实给出局部辛同胚。

	对于给定二次函数，驻点方程为
	\[
		-\frac{Q-q}{t}+\frac{R-Q}{s}=0,
		\qquad Q=\frac{s q+tR}{t+s}.
	\]
	此时 \(H=-(1/t+1/s)I\) 可逆。代回得 \(U(q,R)=-|R-q|^2/(2(t+s))=S_{t+s}(q,R)\)。单步映射为 \((q,p)\mapsto(q+tp,p)\)，所以复合参数相加，也从映射本身验证了生成函数的计算。
\end{solution}

\begin{exercise}\label{g4:ex:0804:rotation-functions}
	对平面旋转
	\[
		(Q,P)=(q\cos\theta-p\sin\theta,\ q\sin\theta+p\cos\theta),
	\]
	分别求以 \((q,Q)\) 和 \((q,P)\) 为变量的生成函数，并确定各自适用的 \(\theta\)。解释为何两种变量选择合起来覆盖全部旋转。
\end{exercise}
\begin{solution}
	记 \(c=\cos\theta\)、\(s=\sin\theta\)。当 \(s\ne0\)，从第一式得 \(p=(cq-Q)/s\)，从第二式得 \(P=(q-cQ)/s\)。因此
	\[
		S_\theta(q,Q)=\frac{c(q^2+Q^2)-2qQ}{2s}
	\]
	满足 \(p=S_q\)、\(P=-S_Q\)，其混合导数为 \(-1/s\)。当 \(c\ne0\)，则 \(p=(P-sq)/c\)、\(Q=(q-sP)/c\)，故
	\[
		F_\theta(q,P)=\frac{qP}{c}-\frac{s}{2c}(q^2+P^2)
	\]
	满足 \(p=F_q\)、\(Q=F_P\)，混合导数为 \(1/c\)。若 \(s=0\)，图像投影到 \((q,Q)\) 的微分退化，因为 \(Q=cq\) 不含 \(p\)；若 \(c=0\)，投影到 \((q,P)\) 同样退化。\(s,c\) 从不同时为零，所以每个旋转至少有一种上述生成函数。公式中的分母记录的是所选投影的失效，而不是辛映射本身的奇异。
\end{solution}

\begin{exercise}\label{g4:ex:0804:local-not-global}
	令 \(q,Q\in\mathbb R^2\)，并设
	\[
		g(Q^1,Q^2)=(e^{Q^1}\cos Q^2,e^{Q^1}\sin Q^2),
		\qquad S(q,Q)=q\cdot g(Q).
	\]
	证明 \(S_{qQ}\) 处处可逆，但方程 \(p=S_q\)、\(P=-S_Q\) 并不定义整个 \(\mathbb R^4_{q,p}\) 上的单值映射。指出局部生成函数定理不能据此推出哪些整体性质。
\end{exercise}
\begin{solution}
	有
	\[
		S_{qQ}=\mathrm{D}g(Q)
		=e^{Q^1}\begin{pmatrix}\cos Q^2&-\sin Q^2\\\sin Q^2&\cos Q^2\end{pmatrix},
		\qquad \det S_{qQ}=e^{2Q^1}>0.
	\]
	然而第一组方程是 \(p=g(Q)\)。当 \(p=0\) 时无解；当 \(p\ne0\) 时，\(Q^1=\log|p|\)，而 \(Q^2\) 可相差任意 \(2\pi\) 的整数倍。因此每个非零 \(p\) 都有无穷多个不同的 \(Q\)，再由 \(P=-\mathrm{D}g(Q)^{\mathsf T}q\) 确定相应 \(P\)。每个选定分支都在足够小的邻域中给出辛同胚，但混合矩阵可逆不保证投影覆盖整个源空间，也不保证全局单射或单值。
\end{solution}

\begin{exercise}\label{g4:ex:0805:projection-area}
	设 \(A\) 是 \(\mathbb R^{2n}_{q^1,p_1,\ldots,q^n,p_n}\) 上的线性辛同构，\(\pi\) 是到第一对坐标平面的投影。证明 \(\pi A(B^{2n}(R))\) 是面积至少为 \(\pi R^2\) 的椭圆。推导这个椭圆的面积公式，从而加强正文的线性非挤压结论。
\end{exercise}
\begin{solution}
	记 \(A\) 的前两行转置为向量 \(a,b\)，并以 \(B\) 表示这两行组成的 \(2\times2n\) 矩阵。令 \(\Omega\) 为配对次序下的标准辛形式矩阵。辛条件给出 \(a^{\mathsf T}\Omega b=1\)，所以两行线性无关，\(BB^{\mathsf T}\) 正定。对给定 \(y\in\mathbb R^2\)，方程 \(Bx=y\) 的最小范数解为
	\[
		x_0=B^{\mathsf T}(BB^{\mathsf T})^{-1}y,
		\qquad |x_0|^2=y^{\mathsf T}(BB^{\mathsf T})^{-1}y.
	\]
	确实，任何其他解与 \(x_0\) 的差属于 \(\ker B\)，而 \(x_0\in\operatorname{im}B^{\mathsf T}=(\ker B)^\perp\)，故范数平方相加。于是投影恰为
	\[
		\{y:y^{\mathsf T}(BB^{\mathsf T})^{-1}y<R^2\},
	\]
	其面积为
	\[
		\pi R^2\sqrt{\det(BB^{\mathsf T})}
		=\pi R^2\sqrt{|a|^2|b|^2-(a\cdot b)^2}.
	\]
	令 \(b_\perp=b-(a\cdot b)a/|a|^2\)。由 \(a^{\mathsf T}\Omega a=0\)、\(\Omega\) 正交及普通 Cauchy--Schwarz 不等式，
	\[
		1=|a^{\mathsf T}\Omega b_\perp|
		\le |a|\,|b_\perp|
		=\sqrt{|a|^2|b|^2-(a\cdot b)^2}.
	\]
	这给出面积下界。如果 \(A(B^{2n}(R))\subset Z^{2n}(r)\)，该椭圆位于面积 \(\pi r^2\) 的圆盘中，便得到 \(R\le r\)。这里还控制了任意形状的线性投影椭圆的面积。
\end{solution}
